% concatenated from main.tex, preamble.tex, intro.tex, envelope.tex, formal.tex, proof.tex, quantum.tex, outro.tex, acks.tex, conic.tex, quantum-proof.tex, interpol-proofs.tex
\documentclass[11pt]{article}
\usepackage[T1]{fontenc}
\usepackage[utf8]{inputenc}
\usepackage{lmodern}
\usepackage[margin=1in]{geometry}
\usepackage{microtype}
\usepackage{amsmath,amssymb,amsthm,mathtools,mathrsfs,bm}
%\usepackage{enumitem}
\usepackage{paralist}
\usepackage{comment}
\usepackage{mdframed}
%\usepackage{textcomp}
%\setlength{\topsep}{0pt}

%\newtheorem{theorem}{Theorem}[section]

\usepackage[dvipsnames]{xcolor}
\usepackage[colorlinks]{hyperref}
\hypersetup{
colorlinks=true,
    linkcolor=red,
    citecolor=OliveGreen,
    filecolor=black,
    urlcolor=black,
}

\allowdisplaybreaks[3]
\numberwithin{equation}{section}


\let\le\leqslant
\let\leq\leqslant
\let\ge\geqslant
\let\geq\geqslant

\let\preceq\preccurlyeq
\let\succeq\succcurlyeq

\usepackage{xfrac}
\usepackage{mathrsfs}
\usepackage{eucal}
\usepackage{dsfont}

\newcommand{\pinv}{\dagger}
\renewcommand{\top}{\mathsf{T}}
\newcommand{\htop}{\mathsf{H}}
\newcommand{\qtop}{\mathsf{Q}}

\newcommand{\e}{{\mathrm{e}}}

\newtheorem{theorem}{Theorem}%[section]
\newtheorem{conjecture}{Conjecture}%[section]
\newtheorem{proposition}[theorem]{Proposition}
\newtheorem{lemma}[theorem]{Lemma}
\newtheorem{corollary}[theorem]{Corollary}
\newtheorem{definition}{Definition}
\theoremstyle{remark}
\newtheorem{remark}[theorem]{Remark}

\newcommand{\odima}[1]{{\color{magenta}#1}}
\newcommand{\droid}[1]{{\color{red}#1}}
\newenvironment{droidblock}{\begingroup\color{red}}{\endgroup}

\newcommand{\proofstep}[1]{$\textbf{#1}^{\circ}$}

\newcommand{\wt}{\widetilde}
\newcommand{\wh}{\widehat}

\newcommand{\dd}{\mathrm{d}}
\newcommand{\R}{\mathbb{R}}
\newcommand{\E}{\mathbb{E}}
\newcommand{\ip}[2]{\left\langle #1,#2\right\rangle}
\newcommand{\norm}[1]{\left\lVert #1\right\rVert}
\newcommand{\ones}{\mathsf{1}_d}
\newcommand{\veps}{\varepsilon}
\newcommand{\sign}{\operatorname{sign}}
\newcommand{\inner}[2]{\langle #1,#2\rangle}
\newcommand{\abs}[1]{\lvert #1\rvert}
\newcommand{\ri}{\operatorname{ri}}
\newcommand{\rbd}{\operatorname{rbd}}
\newcommand{\supp}{\operatorname{supp}}
\newcommand{\conv}{\operatorname{conv}}
\newcommand{\Env}{S}
\newcommand{\Ent}{\textup{\textsf{Ent}}}
% \newcommand{\Pin}{\textup{\textsf{Pin}}}
\newcommand{\Diag}{\textup{Diag}}
\newcommand{\Pin}{\mathsf{\Pi}}
\newcommand{\Orc}{\mathscr{O}}
\newcommand{\Mtd}{\mathscr{M}}
\newcommand{\FOM}{\textsf{FOM}}
\let\Div\relax\newcommand{\Div}{D_h}
\let\Dh\Div
%\newcommand{\Fh}{\mathscr{F}_h}
%\newcommand{\Fh}{\mathcal{F}_h}
\newcommand{\Fh}{\Ent_d}
\newcommand{\Vh}{\textsf{Lip}_d}
\newcommand{\Simp}{\Delta_d}

\newcommand{\cI}{\mathcal{I}}
\newcommand{\ProjTan}{\Pi_{\Tan}}

\newcommand{\bOmega}{\boldsymbol{\Omega}}
\newcommand{\bPhi}{\boldsymbol{\Phi}}

\newcommand{\Herm}{\mathbf{H}^d}
%\newcommand{\Dens}{\mathbf{D}_d}
\newcommand{\Dens}{\boldsymbol{\Delta}_d}
%\newcommand{\Qd}{\mathcal{Q}_d}
\newcommand{\Qd}{\mathsf{Ent}_d^\mathsf{H}}

\newcommand{\qre}[2]{D_{h \hspace{1pt} \circ \lambda}(#1,#2)}
\newcommand{\Tan}{\Theta_d}
\newcommand{\spanof}[1]{\operatorname{span}\{#1\}}
\newcommand{\Opt}{\operatorname{Opt}}
\newcommand{\Risk}{\mathsf{Risk}}
\newcommand{\prox}{\operatorname{prox}}
\DeclareMathOperator*{\argmin}{\arg\!\min}
\DeclareMathOperator*{\argmax}{\arg\!\max}
\DeclareMathOperator*{\Argmin}{\operatorname{Arg}\!\min}
\DeclareMathOperator*{\Argmax}{\operatorname{Arg}\!\max}
\DeclareMathOperator{\tr}{tr}
\DeclareMathOperator{\diag}{diag}


\newcommand{\1}{\mathbf{1}}
\newcommand{\ud}{u_d}
\newcommand{\cl}{\operatorname{cl}}
%\newcommand{\bnabla}{\boldsymbol{\nabla}}
\newcommand{\bnabla}{\nabla}
\newcommand{\onabla}{\overline{\nabla}}
\newcommand{\unabla}{\underline{\nabla}}
\newcommand{\half}{\frac{1}{2}}




\title{Entropy-Smooth Convex Optimization Cannot Be Accelerated}
\author{
Jacob M.~Aguirre\thanks{Georgia Institute of Technology, H.~Milton Stewart School of Industrial and Systems Engineering (ISyE), Atlanta, USA. 
%Supported by the NSF Graduate Research Fellowship under Grant No.~DGE-2039655. 
Email: \texttt{aguirre@gatech.edu}.}
\and 
Dmitrii M.~Ostrovskii\thanks{Georgia Institute of Technology, School of Mathematics \& H.~Milton Stewart School of Industrial and Systems Engineering (ISyE), Atlanta, USA. Email: \texttt{ostrov@gatech.edu}.}
}
\date{}

\begin{document}
\maketitle

\begin{abstract}
We prove an $\Omega(L/T)$ lower bound for the convergence rate of minimization in the class of functions that are convex and~$L$-smooth relative to negative entropy on the standard~$d$-simplex, valid for every first-order method when~$\smash{d = \Omega(T^2)}$. 
In particular, this shows that mirror descent is optimal up to a logarithmic factor in this class. This may be surprising due to the fact that accelerated methods are readily available under the assumption of smoothness in $\smash{\ell_1}$-norm.
While Dragomir et al.~(Mathematical~Programming, 2022) have already showed that acceleration might be impossible under relative smoothness, their prox-function is pathological and constructed together with the hard instance. 
In contrast, we show non-acceleration for a {specific} prox-function with particularly favorable structure.
We also extend the result to the quantum setting, proving the same lower bound in the class of functions~$L$-smooth relative to the negative von Neumann entropy on the spectrahedron of~$d \times d$ Hermitian positive-semidefinite matrices with unit trace. 
%Finally, we describe the tightest interpolant in the class for given interpolation data.
%yet, despite the uniquely favorable combination of Bregman and~$f$-divergence properties, acceleration mechanism breaks down without the norm. 
%This also gives a first explicit example of a pathological potential



%Namely: any algorithm, when given access to the standard (local) first-order oracle that returns the objective value and gradient at a queried point (with~$+\infty$-valued gradient entries allowed on the boundary). 

\end{abstract}

%\medskip
%\noindent\textbf{Keywords}: relative smoothness, entropy geometry, oracle complexity, Bregman--Moreau envelopes, quantum relative entropy
%
%\smallskip
%\noindent\textbf{MSC2020}: 90C25, 90C30, 90C60, 49M37.


\section{Introduction}
\label{sec:intro}

In this paper, we study convex optimization problems on the standard simplex~$\Simp \subset \R^d$, of the form
\begin{equation}
\label{eq:intro-problem}
\min_{x\in \Simp} f(x).
\end{equation}
More precisely, we are interested in the first-order oracle complexity\footnote{In the classical sense of Nemirovski and Yudin~\cite{nemirovski1983problem}; we shall briefly recap their complexity framework in Section~\ref{sec:prelim}.} of the natural problem class in which the objective smoothness is measured with respect to negative Shannon entropy
\begin{equation}
\label{eq:intro-entropy}
h: \Simp \to \R, \qquad h(x) = \sum_{k=1}^d (x)_k \log (x)_k,
\end{equation}
where~$(x)_k$ denotes the~$k{\textup{th}}$ entry of~$x$ and we use the standard convention~$0 \log 0 = 0$.
Here, ``smoothness of one function with respect to another'' refers to the notion of {\em relative smoothness}, as defined in~\cite{BauschkeBolteTeboulle2017} and~\cite{LuFreundNesterov2018},
which we now recall in the form adapted to our setting. 
Any function~$\phi$ strictly convex on a domain~$X \subseteq \R^d$ and $C^1$ in its relative interior, defines the {Bregman divergence}
\[
D_{\phi}: X \times \ri(X) \to \R_+, 
\qquad
D_{\phi}(x,u) =  \phi(x) - \phi(u) - \langle \bnabla\phi(u), x-u \rangle,
\]
a directed measure of discrepancy of an arbitrary point~$x \in X$ from the ``center'' point~$u \in \ri(X)$. 
(In the sequel,~$X$ is assumed to lie in some affine subspace of~$\R^d$, and we write~$\bnabla f(x)$ for the gradient of~$f: X \to \R$ restricted to the affine hull of~$X$; a more detailed discussion is deferred to Section~\ref{sec:prelim}.)
For~$L>0$, a function~$f: X \to \R$ in~$C^1(\ri(X))$ is called~{\em $L$-smooth relative to~$\phi$} if
% one has
%\,\footnote{Utilizing~$\nabla f$ requires some care when~$X$ is not a proper subset of~$\R^d$. These technicalities are deferred to Section~\ref{sec:prelim}.}
\begin{equation}
\label{eq:relative-smoothness}
f(x) - f(u) - \inner{\bnabla f(u)}{x-u} \le LD_{\phi}(x,u) 
\qquad \forall (x,u) \in X \times \ri(X).
\end{equation}
The inequality in~\eqref{eq:relative-smoothness} can be recast as~$D_f(x,u) \le LD_{\phi}(x,u)$, so~\eqref{eq:relative-smoothness} is equivalent to the convexity of~$L\phi - f$.
%~$\nabla^2 f(x) \preccurlyeq L\nabla^2 h(x)$ in all~$C^2$-points.
%but we shall never make use of these alternative forms
When~$X = \R^d$ and~$\phi(\cdot) = \frac{1}{2} \|\cdot\|_2^2$, the associated Bregman divergence is~$\frac{1}{2}\|x-u\|_2^2$, and~$L$-relative smoothness reduces to the usual Euclidean smoothness, i.e.,~$L$-Lipschitzness of~$\bnabla f$.
%
When~$(X,\phi) = (\Simp,h)$, the corresponding Bregman divergence gives the Kullback--Leibler divergence. 
For our purposes, it is convenient to treat the latter as an extended-value function over~$\Simp \times \Simp$,
%with the convention \(0\log0=0\). 
%Since our functions are finite on the closed simplex, it is useful to extend this quantity to all pairs \(x,u\in\Simp\) by setting
\begin{equation}
\label{eq:intro-kl}
        \Div(x,u)
        :=
        \begin{cases}
        \displaystyle
        \sum_{k \,\in\, \supp x} (x)_k\log\frac{(x)_k}{(u)_k},
        & \supp x\subseteq\supp u,\\[1.2ex]
        +\infty,
        & \text{otherwise},
        \end{cases}
\end{equation}
with the convention~$0\log 0 = 0$.
Geometrically,~$\Div(x,u) < +\infty$ occurs when~$u$ is ``at least as interior'' as~$x$;
%the minimal face containing~$x$ is contained in the corresponding minimal face of~$u$.
%Thus \(\Div:\Simp\times\Simp\to[0,+\infty]\) is an extended-valued divergence.  
%It is finite exactly when the second argument is at least as interior as the first, in the sense that \(\supp x\subseteq\supp u\).  
on \(\Simp\times\ri(\Simp)\) this definition coincides with the usual Bregman divergence generated by~$h(\cdot)$. 

%Another classic setting is the one with~$X = \Simp$,~$\phi = h$, and the Kullback-Leibler (KL) divergence
%%\begin{equation}
%%\label{eq:intro-kl}
%%\Dh(x,u)=\sum_{i=1}^d x_i \log \frac{x_i}{u_i},
%%\end{equation}
%with the usual convention~$0 \log 0 = 0$, as the Bregman divergence.
%Note that~$\Dh(x,u) < +\infty$ for all~$x \in \Simp$ if~$u \in \ri(\Simp)$; more generally:~$\Dh(x,u) < +\infty$ when~$u \gg x$ as distributions, i.e.~when~$\supp(u) \supseteq \supp(x)$;
%~$\supp(u) \supseteq \supp(x)$, where~$\supp(\cdot)$ is the set of nonzero entries of a vector in~$\R^d$. 
%(Geometrically, this means that for~$u$ on some~$k$-dimensional face of~$\Simp$, the ``center''~$v$ must be in the relative interior of that face.)
%
We are interested in the~$(\Simp,h)$ case, so let us formally define the corresponding class of functions.

\begin{definition}
\label{def:entropy-smoothness}
The class~$\Fh(L)$ consists of all functions~$f: \Simp \to \R$ that are convex, continuously differentiable on~$\ri(\Simp)$, and satisfy the following inequalities for all~$(x,u) \in \Simp \times \ri(\Simp)$:
%~$L$-smooth relative to~$h$ on~$\Simp$; in other words,
\begin{equation} 
\label{eq:intro-sandwich} 
0 \le f(x)-f(u)-\inner{\bnabla f(u)}{x-u} \le L\Div(x,u).
\end{equation}
%for every~$x \in \ri(\Simp)$ and~$y \in \Simp$.
\end{definition}
%
\noindent
The first inequality in \eqref{eq:intro-sandwich} is the usual first-order convexity certificate, whereas the second one is the condition of~$L$-smoothness relative to~$h$ on~$\Simp$.
%the entropy relative-smoothness condition. 
Note that differentiability is only required on \(\ri(\Simp)\); yet, functions in \(\Fh(L)\) are finite and continuous on the closed simplex. In particular,~$h\in\Fh(1)$.

\iffalse
Assuming~$f \in C^1(\ri(\Simp))$, the inclusion~$f \in \Fh(L)$ amounts to convexity of~$f$ and~$Lh - f$; explicitly,
\begin{equation} 
\label{eq:intro-sandwich} 
0 \le f(y)-f(x)-\inner{\nabla f(x)}{y-x} \le L\Div(y,x).
\end{equation}
where the right-hand side is the KL divergence of~$y$ from~$x$, in that order (cf.~\eqref{eq:intro-kl}). 
We shall be referring to functions satisfying the right inequality in~\eqref{eq:intro-sandwich} as {\em $L$-entropy smooth}. 
Note that Definition~\ref{def:entropy-smoothness} does not require~$f$ to be differentiable on the relative boundary of~$\Simp$; e.g.,~$h \in \Fh(1)$.
\fi

The notion of relative smoothness, proposed simultaneously by~Bauschke et al.~\cite{BauschkeBolteTeboulle2017} and Lu et al.~\cite{LuFreundNesterov2018}, is motivated by the simple yet striking observation: inequality~\eqref{eq:relative-smoothness} 
%(with arbitrary~$X$ and~$\phi$ instead of~$\Delta$ and~$h$) 
can be interpreted as the proximal descent lemma formulated directly in terms of the Bregman divergence at hand.
%without using any norm. 
%Indeed, recall 
%\footnote{Or distance-generating function, in the terminology of A.~Nemirovski~\cite{ben2001lectures}. We use the term ``potential'' for brevity.} 
Using this observation,~\cite{BauschkeBolteTeboulle2017} and~\cite{LuFreundNesterov2018} generalized the classical analysis of mirror descent~\cite{nemirovski1983problem}, showing that
%a.k.a.~the~Bregman gradient method, leading to the bound
\begin{equation}
\label{eq:intro-upper-general}
f(x_T)- \min_{x \in X} f(x)\le \frac{LD_{\phi}(x^\star,x_0)}{T} \qquad \forall x^\star \in \Argmin\limits_{x \in X} f(x)
\end{equation}
after~$T$ iterations of the algorithm initialized from~$x_0$,  all without using {\em any norm} on the domain.
%where
%\[
%\Omega_{X}(\phi) := \max_{u \in X} \phi(u) - \min_{v \in X} \phi(v) 
%\]
%is the variation of~$\phi$ on~$X$.x
Specializing this to the entropy on~$\Simp$ with initialization at the barycenter~$d^{-1} \ones$, we get the bound
\begin{equation}
\label{eq:intro-upper-KL}
f(x_T)- \min_{x \in \Simp} f(x)\le \frac{L\log d}{T} \qquad \forall f \in \Fh(L).
\end{equation}
%proved without any appeal to the~$\ell_1$-norm or Pinsker's inequality. 
To put this into perspective, the classic analysis of mirror descent 
%(e.g.~\cite{juditsky2011first}) 
proceeds by choosing a norm~$\|\cdot\|$ such that~$\phi$ is~$1$-strongly convex w.r.t.~$\|\cdot\|$ on~$X$, 
%granting convexity and~$L$-smoothness in~$\|\cdot\|$, i.e.
which implies convexity and~$L$-smoothness in~$\|\cdot\|$,
%assuming~$L$-smoothness in the~$\ell_1$-norm, i.e.~the inequalities 
\begin{equation*}
%\label{eq:intro-norm-smoothness}
0 \le f(x) - f(u) - \inner{\bnabla f(u)}{x-u} \le \frac{L}{2} \|x-u\|^2
\qquad \forall x,u \in X,
\end{equation*}
and using~$1$-strong convexity of~$\phi$ at some stage of the analysis. 
This route still results in~\eqref{eq:intro-upper-general}, but only in the smaller class of functions: indeed, the above inequality implies~\eqref{eq:relative-smoothness} but not vice versa.
In particular, in the relevant to us case~$(\Simp,h)$, the suitable norm is~$\|x\|_{1} = \sum_{i = 1}^d |(x)_i|$;
%and the~$1$-strong convexity of~$h$ in~$\ell_1$-norm is Pinsker's inequality; 
adopting the name~$\Vh(L)$ for the corresponding class of functions on~$\Simp$, i.e.~those satisfying the inequalities
\begin{equation}
\label{eq:intro-l1-smoothness}
0 \le f(x) - f(u) - \inner{\bnabla f(u)}{x-u} \le \frac{L}{2} \|x-u\|_1^2
\qquad \forall x,u \in \Simp,
\end{equation}
the inclusion~$\Vh(L) \subsetneq \Fh(L)$ follows from  the~$1$-strong convexity of~$h$ w.r.t.~$\ell_1$-norm, that is Pinsker's inequality~\cite{pinsker1964information}.  In these terms, the guarantee in~\eqref{eq:intro-upper-KL} extends from~$\Vh(L)$ to~$\Fh(L)$.
\iffalse
that are stronger than~\eqref{eq:intro-sandwich} on the right, and then using Pinsker's inequality~$\frac{1}{2}\|u-v\|_1^2 \le \Dh(u,v)$, a.k.a.~$1$-strong convexity of~$h$ on~$\Simp$ in~$\ell_1$-norm, at some stage of the analysis. 
While this results in the same error bound as in~\eqref{eq:intro-upper-KL}, the class~$\Vh(L)$ of functions satisfying~\eqref{eq:intro-l1-smoothness} is genuinely smaller:
\fi

One may ask whether other classical results for proximal algorithms admit similar generalizations. 
In particular, a very natural question, mentioned already in~\cite{BauschkeBolteTeboulle2017},~\cite{LuFreundNesterov2018} and more explicitly discussed by R.-A.~Dragomir in his PhD thesis~\cite{DragomirThesis2021}, is whether acceleration ``\`a la Nesterov''~\cite{nesterov1983method} is possible under relative smoothness.  
Indeed, in the class of functions convex and~$L$-smooth on a set~$X$ with respect to a given norm~$\|\cdot\|$, 
%generalizing~\eqref{eq:intro-l1-smoothness}, 
one can obtain~$O(LD_{\phi}(x^\star,x_0)T^{-2})$ error using any potential that is~$1$-strongly convex with respect to the norm~$\|\cdot\|$; see~\cite{nesterov2013first}.\,\footnote{We write~$g = O(f)$ or~$f = \Omega(g)$ if there exists a universal constant~$c > 0$ such that~$g \le cf$ for all argument values.}
With this in mind, we are led to the question:
%\begin{mdframed}
\begin{quote}
{\em
Given~a convex domain~$X \subseteq \R^d$ and a potential~$\phi$,  is there a first-order algorithm with guaranteed~$O(T^{-2})$ convergence in the corresponding class of relatively smooth functions?
\em}
\end{quote}
%\end{mdframed}
One may also specialize this question to {\em concrete} domain-potential pairs. 
In particular, for~$(\Simp,h)$
%the known guarantee is
\begin{equation}
\label{eq:intro-upper-KL-acc}
f(x_T)- \min_{x \in \Simp} f(x) = O\left(\frac{L\log d}{T^2}\right) \qquad \forall f \in \Vh(L),
\end{equation}
and the question we are left with is the following one:
\begin{mdframed}
\begin{quote}
\begin{center}
{\em
Can the guarantee in~\eqref{eq:intro-upper-KL-acc} be extended to the class~$\Fh(L)$?
\em}
\end{center}
\end{quote}
\end{mdframed}
%In particular, in the case of~$(\Simp, h)$, 
%in particular,
%that is~$1$-strongly convex w.r.t.~$\|\cdot\|$.
%class~$\Vh(L)$ admits iterative methods with stronger convergence guarantees (see~\cite{nesterov2013first}):
%in other words, one may ask for extending the property holding 
%where~$O(\cdot)$ hides a universal constant.
To the best of our knowledge, both these questions are open.
Our paper resolves the second question.

\paragraph{Our results.}
Focusing on the case of~$(\Simp,h)$, we answer the acceleration question in the negative.
%\odima{
A rigorous statement of our result relies on the basic notions of black-box complexity theory~\cite{nemirovski1983problem}, to be recapped in Section~\ref{sec:prelim}.
%\ref{sec:formal}.
The simplified formulation, presented next, suffices to make our point.
%}
\begin{theorem}
\label{th:intro-lower}
Let~$L > 0$,~$T \ge 1$, and~$d = \Omega(T^2)$.
%deterministic local first-order method that
For any deterministic method that makes~$T$ queries~$x_{1}, \dots, x_{T} \in \ri(\Simp)$ of the oracle~$x \mapsto (f(x),\bnabla f(x))$ and returns~$x_{T+1} \in \Simp$, there exists~$f(\cdot) \in \Fh(L)$ such that
\begin{equation}
\label{eq:intro-lower}
f(x_{T+1}) - \min_{x \in \Simp} f(x) > \frac{L}{4(T+1)}. 
\end{equation}
\end{theorem}
%
This result shows that, in contrast to~$\Vh(L)$, accelerated~$O(T^{-2})$ convergence cannot be attained on~$\Fh(L)$. 
Moreover, the convergence guarantee~\eqref{eq:intro-upper-KL} of entropic mirror descent with stepsize~$1/L$ is optimal---up to a logarithmic factor---if the dimension is large enough, namely when~$\smash{d =\Omega(T^2)}$.

Let us make several remarks regarding Theorem~\ref{th:intro-lower}. 
%deferring further discussion to Section~\ref{sec:outro}. 
%
%\vspace{-0.1cm}
\begin{itemize}
%\begin{compactenum}
%\vspace{0.1cm}
\item[R1.] As mentioned, Theorem~\ref{th:intro-lower} leaves a logarithmic gap between the best available upper and lower bounds for the class~$\Fh(L)$; cf.~\eqref{eq:intro-upper-KL}. 
%\odima{Closing this gap appears to be a difficult problem; a possible route is outlined in Section~\ref{sec:outro}.}
We expect that the lower bound in~\eqref{eq:intro-lower} is loose, and the missing~$\log d$ factor can be recovered.
A promising approach is outlined in Section~\ref{sec:outro}.
%and implement the first step in this direction.

%\vspace{0.1cm}
\item[R2.] When~$d = \exp(T)$, the right-hand side of~\eqref{eq:intro-lower} reads as~$O(\smash{{LT^{-2} \log d}})$, matching the upper bound in~\eqref{eq:intro-upper-KL-acc} and diverging from~\eqref{eq:intro-upper-KL} by~$T$.  Therefore, acceleration is formally not ruled out for such ``extremely high-dimensional'' problems; however, this regime is of limited practical interest anyway, since even a single iteration of a deterministic first-order method is prohibitive.
%However, in this regime the guarantees involving~$\log d$ are hardly satisfactory anyway; 
%accordingly, proximal methods suffering from this overhead are replaced with conditional gradient methods.

%\vspace{0.1cm}
\item[R3.] The quadratic dependence of the dimension on~$T$  is crucial for our resisting oracle construction.
%in the proof of Theorem~\ref{th:intro-lower}. 
This gives an~$O(T^{-2})$ ``safety margin'' that ensures consistency of the transcript after~$T$ steps.

%\vspace{0.1cm}
\item[R4.] Restriction to interior queries is natural: indeed,~$f$ might be non-differentiable on the relative boundary; this is not a pathological case either, as shown by the scaled negative entropy~$Lh$.
%A possible alternative path is to allow to query the oracle on the boundary, and let it return finite 
%Boundary queries can be emulated via a compactness argument in Appendix~\odima{\bf ?}
%A method may emulate boundary queries from the interior; tje
%namely as limits of oracle queries at points approaching the boundary from \(\ri(\Simp)\).
\vspace{-0.1cm}
%Alternatively, we could extend the definition of the first-order
%We formulate the first-order oracle for interior queries, where the differential is an ordinary affine gradient.  Boundary queries may be incorporated by the standard closure convention, namely as limits of oracle calls at points approaching the boundary from \(\ri(\Simp)\).
%\end{compactenum}
\end{itemize}

Theorem~\ref{th:intro-lower} admits a noncommutative generalization, pertaining to functions  on the ``spectraplex''
\[
\Dens := \{X \in \Herm_+: \tr(X) = 1\}
\] 
where~$\Herm_+$ is the Hermitian positive-semidefinite cone. 
In this setting, defined rigorously in Section~\ref{sec:quantum}, 
we introduce the class~$\Qd(L)$ of functions on~$\Dens$ that are convex,~$C^1(\ri(\Dens))$, and~$L$-smooth relative to the negative von Neumann entropy~$\tr(X \log X)$, which is the spectral analog of~$h(x)$. 
It turns out that this noncommutative setting can be reduced to the diagonal case, corresponding to the class~$\Fh(L)$ and the setting of Theorem~\ref{th:lower}. 
%(spectrahedron of unit-trace positive-semidefinite Hermitian matrices) 
This is done by embedding the hard instance of Theorem~\ref{th:intro-lower} diagonally, and using Lindblad's inequality~\cite{Lindblad1975} to argue that the composition~$F(\cdot) = f( \diag(\cdot))$ of~$f(\cdot) \in \Fh(L)$ and the diagonal extraction map belongs to~$\Qd(L)$. The final ingredient is the observation that the diagonal of the (matrix) transcript of an arbitrary first-order method run on~$\smash{F}$ emulates the transcript of some first-order method run on~$f$.
We defer further details to Section~\ref{sec:quantum}.

Finally, as a minor contribution, in Section~\ref{sec:outro} we find an explicit form of the pointwise minimal interpolant in the class~$\Fh(L)$ for given interpolation data~$\cI = (x_i,f_i,\xi_i)_{i = 1}^{N}$, defined as the function~$f \in \Fh(L)$ that interpolates~$\cI$ to first order, i.e.~satisfies~$f(x_i) = f_i$ and~$\nabla f(x_i) = \xi_i$ for all~$i \in [N]$, and is no larger than any other such function at every point~$x \in \Simp$. This result follows easily from the results in R.-A.~Dragomir's PhD thesis~\cite{DragomirThesis2021}; we record it due to its pivotal role in the promising approach of improving Theorem~\ref{th:intro-lower}. Further details are deferred to Section~\ref{sec:outro}.

\paragraph{Summary of the approach.}
Our construction is based on the {\em right Bregman--Moreau envelope}~\cite{BauschkeDaoLindstrom2018} applied to an incremental coordinatewise construction in the spirit of~Guzm\'an and Nemirovski~\cite{GuzmanNemirovski2015}.
The right Bregman-Moreau envelope replaces the inf-convolution smoothing operator of~\cite{GuzmanNemirovski2015} as the smoothing mechanism.
This smoothing mechanism heavily relies upon the $f$-divergence properties of the KL divergence, namely its joint and separate convexity in {both} arguments;
as demonstrated in~\cite{BauschkeBorwein2001}, among Bregman divergences, these properties occur only in
the Euclidean and KL cases.
%so its employment in our context is natural.
%This smoothing mechanism is well-adapted to the entropy setting, 
Crucially, the right envelope of a convex function is convex and smooth with respect to the potential~(see~\cite{BauschkeDaoLindstrom2018});
meanwhile, the local smoothing of~\cite{GuzmanNemirovski2015} cannot be employed, as it would give an instance in~$\Vh(L)$, and the latter class {\em does} admit acceleration. 
%(cf.~\eqref{eq:intro-upper-KL-acc}). 
%While the smoothing property holds in the general Bregman case, 
%In fact, this envelope is correctly defined 
%Moreover, joint convexity of~$\Div(u,v)$ gives {\em convexity} of the envelope (by the preservation of convexity under partial minimization), as per~\eqref{eq:intro-sandwich}.

%while we aim for~$\Fh(L) \setminus \Vh(L)$.

%This is natural since indeed,~$S_{\eta}[\cdot]$ operates directly with the Bregman structure, and in our~$(\Simp, h)$ setting one may leverage the joint convexity of~$\Div(u,v)$ 

%serves as the direct entropy-geometry counterpart of the local smoothing in \cite{GuzmanNemirovski2015}. In the entropic setting the relevant locality neighborhood is
%\begin{equation}\label{eq:intro-neighborhood}
%N_{\eta,G}(x):=\{z\in \Simp^\circ:\|x-z\|_z\le \eta G\},
%\end{equation}
%where local norms are defined on the affine hull of the simplex. The prox point satisfies the exact affine displacement identity
%\begin{equation}\label{eq:intro-displacement}
%\|x-u_\eta[g](x)\|_{u_\eta[g](x)}
%=
%\eta \|v\|_{u_\eta[g](x),*},
%\qquad
%v\in \partial g(u_\eta[g](x))
%\end{equation}
%for the subgradient appearing in the optimality system. The smoothing loss is then controlled not by a general one-sided convexity inequality, but by atom variation on the induced locality set. For max-affine atoms this gives the bound $\eta G^2$; for the balanced logarithmic atoms it gives the bound $8\eta$. This atomwise locality allows for the two base functions that agree on a neighborhood containing the relevant prox points yield identical shifted envelopes and identical local-oracle answers.

%We consider two atom families. The first is linear: maxima of affine functionals built from the hypercube sign matrix already yield an $\Omega(L/T)$ lower bound. The second is logarithmic. In that case the coordinate atoms $-\log x_i$ satisfy the exact identity $\|\nabla(-\log x_i)\|_{x,*}^2=(1-x_i)/x_i$, which forces every active coordinate to satisfy $x_i\gtrsim\eta^2$ and therefore rules out singleton-support constructions. Balanced half-hypercube supports remove this obstacle and again yield an $\Omega(L/T)$ lower bound.

%\vspace{0.1cm}
%Let is briefly discuss the additional related work, then give the roadmap for the rest of the paper.

\paragraph{Related work.}
%
%The relative-smoothness framework, controlling $f$ through the convexity of $Lh-f$, originates in entropic proximal geometry and was developed systematically in \cite{BauschkeBolteTeboulle2017} and~\cite{LuFreundNesterov2018}. The Bregman--Moreau envelope machinery and its differential properties were treated abstractly in \cite{BauschkeDaoLindstrom2018}; Chapter 5 of \cite{DragomirThesis2021} specializes performance estimation and interpolability to the entropy kernel. Dragomir, Taylor, d'Aspremont, and Bolte \cite{DragomirTaylordAspremontBolte2022} show that the \emph{generic} relatively smooth class does not admit a universal accelerated rate. Their result, however, leaves open the question for the fixed entropy kernel on the simplex; this is precisely the gap we address.
%
In addition to~\cite{GuzmanNemirovski2015},  closely relevant to ours is the work of Dragomir, Taylor, d'Aspremont, and Bolte~\cite{DragomirTaylordAspremontBolte2022}, who showed the following: 
there {\em exists a pair}~$(f,\phi)$ in which~$\phi$ is strictly convex on the positive orthant~$\R^d_+$,~$f$ is convex and smooth relative to~$\phi$,
%in the sense of~\eqref{eq:relative-smoothness}, 
and any first-order method has worst-case convergence rate no better than~$\Omega(1/T)$. 
Their result is based on so-called performance estimation techniques (e.g.~\cite{DroriTeboulle2014},~\cite{TaylorHendrickxGlineur2017}), allowing one to find worst-case instances in infinite-dimensional functional classes by reducing the corresponding optimization problem to a (finite-dimensional) semidefinite program (SDP), dubbed ``performance estimation program'' (PEP).
Usually, the PEP methodology is limited to Euclidean geometry, since the SDP representation arises from the Gram matrix that juxtaposes the candidate iterates~$x_{k}$ and gradients~$g_{k} = \nabla f(x_{k})$ in the transcript; as a result, PEPs are poorly suited for dealing with non-Euclidean geometries, where the terms of the form~$\| x_k - x_l \|^2$ and~$\|g_k - g_l\|_*^2$ cannot be expressed linearly via dot products. 
The authors of~\cite{DragomirTaylordAspremontBolte2022} elegantly sidestepped this limitation by allowing the potential itself to vary with~$f$, so that the PEP constructs a worst-case {\em pair}~$(f,\phi)$. 
%
However, the resulting potential is quite pathological---as one might expect with its adversarial origin---and it may very well be that some {\em specific} pairs~$(X,\phi)$ do admit acceleration. 
%moreover, the methodology of selecting~$\phi$ adversarially is in logical contradiction with both the philosophy and the established practice of picking~$\phi$ favorable to the feasible set. 
While our result shows this is not the case for~$(\Simp, h)$, other highly structured settings remain open, most notably that of the logarithmic-barrier potential~$\phi(x) = -\sum_{i = 1}^d \log(x)_i$ on~$\R^d_+$. 


%in fact, this realization motivated the present 
%In particular,~$(\Simp, h)$ seems to be a perfect candidate, due to the uniquely favorable algebro-analytic structure of~$h$, 
%Indeed, as shown by Bauschke and Borwein~\cite{BauschkeBorwein2001}, KL divergence is one of the only two {\em jointly convex} Bregman divergences, the other one being~$\frac{1}{2} \|u-v\|_2^2$ (plus derived from these via ``natural algebraic operations;" see~\cite{BauschkeBorwein2001} for more details). This leads to {exact} interpolation conditions, not available in the general Bregman scenario~\cite[Chap.~5 and Sec.~5.3.3]{DragomirThesis2021}. 

%\odima{Write about reverse information projection?}

%On a separate but closely related note, 
Chapter~5 of Dragomir's PhD thesis~\cite{DragomirThesis2021} provides {exact}  first-order interpolation conditions for the class~$\smash{\Fh^+(L)}$ of functions on the nonnegative orthant~$\smash{\R^d_+}$ that are convex and smooth relative to the unnormalized entropy~$\smash{h_\e(x) := \sum_{k = 1}^d (x)_k \log (x)_k - (x)_k}$; such conditions are the inequalities imposed on the interpolation data~$\cI = \smash{(x_i,f_i,\xi_i)_{i = 1}^N}$, whose validity is equivalent to the existence of a function in~$\Fh^+(L)$ that interpolates~$\cI$. Conceptually, these conditions are counterparts of SDP-representable interpolation conditions in the Euclidean case, which are the crux of the PEP framework, and their existence crucially relies upon the joint convexity of KL divergence~(see \cite[Lem.~5.3.3 and Rem.~3]{DragomirThesis2021}). In Section~\ref{sec:outro}, we generalize these conditions for the class~$\Fh(L)$ of entropy-smooth functions on~$\Simp$, and then use them to derive the {pointwise minimal} interpolant.
 It appears that these results might be leveraged to 
%improve Theorem~\ref{th:intro-lower} and 
recover the logarithmic factor missing in~\eqref{eq:intro-lower}. In Section~\ref{sec:outro} we further discuss this possibility in the light of the recent work of Florea and Nesterov~\cite{florea2025optimal}.

Finally, we mention the triangular scaling exponent framework of~\cite{HanzelyRichtarikXiao2021}, which seeks to relax the standard assumption of~$1$-strong convexity of~$\phi$ w.r.t.~a norm while retaining accelerated convergence. 
Our critique is that, while leading to locally adaptive algorithms that prove to be highly effective for optimization problems arising in some applications, the triangular scaling condition is hard to ensure in the worst case.
In particular, Theorem~\ref{th:intro-lower} shows that this condition is not satisfied for~$\Fh(L)$.
%and one may directly verify that it fails in the logarithmic-barrier setting.
%In conclusion, we mention the recent works~\cite{FanKroerFarina2024,KornowskiShamir2025,SchlisselbergShermanKorenMansour2025} 
%obtain online-learning lower bounds for dilated entropy on treeplexes; Kornowski and Shamir \cite{KornowskiShamir2025} treat simplex-based matrix games under matrix-access oracles; Schlisselberg, Sherman, Koren, and Mansour \cite{SchlisselbergShermanKorenMansour2025}
%In particular, none of these results establishes a direct local-oracle lower bound for $\Fh(L)$ on $\Simp$. 
% target ~$\Fh(L)$ setting; in particular, based  In adjacent models, Guzm\'an and Nemirovski \cite{GuzmanNemirovski2015} derive smooth local-oracle lower bounds over $\ell_p$-balls via Moreau-type local smoothing and, for $1\le p<2$, through nearly Euclidean sections. 
% Fan, Kroer, and Farina \cite{FanKroerFarina2024} obtain online-learning lower bounds for dilated entropy on treeplexes; Kornowski and Shamir \cite{KornowskiShamir2025} treat simplex-based matrix games under matrix-access oracles; Schlisselberg, Sherman, Koren, and Mansour \cite{SchlisselbergShermanKorenMansour2025} identify a separate entropy-specific barrier for approximate online mirror descent. 


\paragraph{Roadmap.}
In Section~\ref{sec:prelim} we collect the ingredients for the proof of Theorem~\ref{th:intro-lower}.
To that end, we recall the properties of the right Bregman-Moreau envelope and the associated smoothing operator, formalize the oracle complexity framework, and describe the family of hard instances subsequently used in the proof.
In Section~\ref{sec:proof} we carry out the proof of Theorem~\ref{th:intro-lower}, and in Section~\ref{sec:quantum} we formulate and discuss its noncommutative generalization. 
In Section~\ref{sec:outro}, we derive exact interpolation conditions and the form of pointwise minimal interpolant in~$\Fh(L)$.


\paragraph{Notation.}
We denote~$[d] := \{1,2,\dots, d\}$. We let~$e_i$ be the~$i$th canonical basis vector in~$\R^d$ and use the concise notation~$(x)_i := \langle x, e_i \rangle$ for the~$i$th entry of~$x \in \R^d$. We let~$\ones$ be the all-ones vector in~$\R^d$. As previously mentioned, we write~$\bnabla f$ for the {tangent} gradient of~$f \in C^1(\ri(\Simp))$; in particular,~$\bnabla f(x) \in \spanof{\ones}^\perp$ for any~$x \in \ri(\Simp)$. 
Additional notation is introduced as necessary.



% envelope.tex -- canonical current version
% cache-stamp: 2026-06-22T07:32:00Z

\section{Building blocks}
\label{sec:prelim}

\paragraph{Tangent gradients.}
In Sections~\ref{sec:prelim}--\ref{sec:proof}, we work with functions defined on~$\Simp$ and differentiable in~$\ri(\Simp)$; this includes every~$f \in \Fh(L)$ and, in particular, the hard instances presented in Section~\ref{sec:hard-subclass}. 
While such functions may arise as restrictions of~$C^1(\R_{++}^d)$-functions (as, for example,~$h$ itself), this is not required in general. For such a function~$f$, we might differentiate it at~$x \in \ri(\Simp)$ intrinsically over the affine hull of~$\Simp$, by identifying the affine subspace~$\spanof{\ones}^{\perp} + d^{-1} \ones$  with~$\R^{d-1}$, via an affine homeomorphism (with orthogonal linear transformation), taking the full gradient of the resulting function in~$\R^{d-1}$, and changing the basis to account for the transformation. 
We take this as the definition of the tangent gradient~$\bnabla f$. 
Note that~$\bnabla f$ always belongs to the tangent space~$\Tan := \spanof{\ones}^\perp$; moreover, if~$f$ is defined in a proper~$\R^d$-neighborhood of~$x$, then~$\bnabla f(x)$ coincides with the Euclidean projection onto~$\Tan$ of the full gradient; denoting the latter with~$\nabla_\star f$, 
\begin{equation}
\label{eq:grad-tangent}
\bnabla f(x) := \nabla_\star f(x) - d^{-1}  \langle \nabla_\star f(x), \ones\rangle \ones, \qquad x \in \ri(\Simp).
\end{equation}
In fact, even if~$f$ is defined only over~$\Simp$, this formula remains valid if we extend~$f$ to~$\R^d_+$ with differentiability over~$\R^d_{++}$. Such an extension is not unique, and~$\nabla_* f$ in general depends on the chosen extension; however,  the dependence is only in the normal component~$d^{-1}\langle \nabla_\star f(x), \ones \rangle \ones$, whereas~$\bnabla f(x)$ is an invariant depending only on the values of~$f$ in a~$\Tan$-neighborhood of~$x \in \ri(\Simp)$.

%Note also that, in contrast with~$\nabla_\star f(x)$, the tangent gradient~$\bnabla f(x)$ is  in this case defined by the  as a result, any extension of~$f$ 



%For a function~$f$ on~$\Simp$, we denote with~$\boldsymbol{\nabla}f$.  
\subsection{Smoothing with right Bregman--Moreau envelope}
\label{sec:envelope}

%In this section, we recall the elementary properties of the right Bregman--Moreau envelope that we shall exploit.
Let $g: \Simp \to \mathbb{R}$ be convex and continuous.
Adapting the standard definition to our setting, {\em the right Bregman--Moreau envelope} of~$g$ with parameter~$\eta > 0$ is defined by its values on the simplex:
\begin{equation}
\label{eq:envelope}
\Env_\eta[g](x) := \min_{u \in \Simp}
\left\{g(u)+\eta^{-1} \Dh(x,u)\right\},
\qquad x \in \Simp.
\end{equation}
For~$x \in \ri(\Simp)$, the minimizer is unique and attained on~$\ri(\Simp)$; 
this defines~$\prox_{\eta g}: \ri(\Simp) \to \ri(\Simp)$,
%~$\prox_{\eta g}: \ri(\Simp) \to \ri(\Simp)$
\begin{equation}
\label{eq:prox}
\prox_{\eta g}(x) = \argmin_{u \in \Simp} \left\{ \eta g(u) + \Div(x,u) \right\}, \qquad x \in \ri(\Simp),
\end{equation}
called {\em the prox-mapping} of $\eta g(\cdot)$.
%To make the paper self-contained, we collect the properties of these objects in subsequent lemmas.
%;~$p_g(\cdot)$ is called the proximal mapping. 
%%Note that~$\Env_\eta[g]$ can be continuously extended from~$\Simp$ to~$\R^d_{+}$, with differentiability in~$\R^d_{++}$, by adding the normalization term~$\eta^{-1} (1-\langle \ones, x \rangle)$ under the minimum, so that~$\Dh(x,u)$ gets augmented to the unnormalized (extended-value)~KL divergence~from~$x \in \R_+^d$ to~$u \in \Simp$.
%In principle, this allows to compute~$\bnabla \Env_\eta[g]$ using~\eqref{eq:grad-exterior}, but we shall obtain it directly. 
%
%\begin{equation}
%\label{eq:right-bm-envelope}
%    \Env_\eta[g](x)
%    :=
%    \min_{u\in\Simp}
%    \big\{
%        g(u)+\eta^{-1} \Div(x,u)
%    \big\},
%    \qquad x\in\Simp.
%\end{equation}
%Note also that minimization in~\eqref{eq:envelope} is effectively restricted to the relative interior of the lowest-dimensional face of~$\Simp$ containing~$u$; in particular, for any~$x \in \ri(\Simp)$ the minimizer stays in~$\ri(\Simp)$. 
%When $x\in\ri(\Simp)$, we denote the minimizer by $p_g(x)$.

Next, we collect some properties of the right Bregman--Moreau envelope and proximal mapping.

\begin{lemma}
%[Envelope regularity]
\label{lem:prox-mapping}
For~$\eta > 0$, the prox-mapping~$\prox_{\eta g}(\cdot)$ is well-defined and continuous on~$\ri(\Simp)$.
%Moreover, the image of any non-vertex~$x \in \Simp$ is supported in the relative interior of the lowest-dimensional face of~$\Simp$ containing~$x$.
%Moreover, 
%For every $x\in\ri(\Simp)$, the minimum in \eqref{eq:right-bm-envelope} is attained at a unique point $p_g(x)\in\ri(\Simp)$.  The map $x\mapsto p_g(x)$ is continuous on $\ri(\Simp)$.  
\end{lemma}

\begin{proof}
As a continuous function on~$\Simp$,~$g$ is bounded from below.
From this and the expression~\eqref{eq:intro-kl} for KL divergence, we see that for~$x \in \ri(\Simp)$, minimization in~\eqref{eq:envelope} can be restricted to~$\ri(\Simp)$.
The same expression shows that the Hessian of~$\Div(x,\cdot)$ is diagonal with positive entries over~$\ri(\Simp)$, so the objective in~\eqref{eq:envelope} is strictly convex and the minimizer is unique.\footnote{Note that~$\Div(x,\cdot)$ is infinitely differentiable as a function over~$\R^d_{++}$, hence we may consider its full Hessian here.}
Continuity of~$\prox_{\eta g}(\cdot)$ follows from the optimal-set mapping theorem of Rockafellar and Wets (see~\cite[Example~5.22]{rockafellar1998variational}), applied to
\[
        \Phi(u,x)
        :=
        g(u)+\eta^{-1}\Div(x,u)+\delta_{\Simp}(u),
        \qquad x\in\ri(\Simp).
\]
Since~$\Simp$ is compact, the function $\Phi$ is proper, lower semicontinuous, and level-bounded in $u$ locally uniformly in $x$.
Since $\Phi(\prox_{\eta g}(\bar x),x)$ is continuous in $x$ at $\bar x$, the cited theorem implies continuity of the value function at any~$\bar x\in\ri(\Simp)$, and outer continuity of the corresponding minimizing set-mapping.
Since this mapping outputs a singleton, this is the continuity of~$x \mapsto \prox_{\eta g}(x)$.
%In our case, this is a singleton, this outer semicontinuity is exactly the continuity of~$p_g(\cdot)$ at $\bar x.$
\end{proof}

\begin{lemma}
\label{lem:envelope-regularity}
The function~$\Env_\eta[g](\cdot)$ is convex on~$\Simp$ and continuously differentiable on~$\ri(\Simp)$, with
\begin{equation}
\label{eq:right-bm-prox-identity}
    \bnabla \Env_\eta[g](x)
    =
    -\eta^{-1}(\log\prox_{\eta g}(x) - \log x),
%    + \lambda(x) \ones
\end{equation}
where~$\log(\cdot)$ is applied entrywise. 
Moreover,~$\Env_\eta[g](\cdot)$ satisfies the inequality for $x \in \ri(\Simp)$ and~$y \in \Simp$:
\begin{equation}
\label{eq:right-bm-relative-smoothness}
    \Env_\eta[g](y)
    \le
    \Env_\eta[g](x)
    +\left\langle \nabla \Env_\eta[g](x),y-x\right\rangle
    +\eta^{-1}\Div(y,x).
\end{equation}
In particular,~$\Env_\eta[g](\cdot) \in \Ent(\eta^{-1})$. 
\end{lemma}

\begin{proof}
By the perspective rule applied coordinatewise to~$x \mapsto x \log x$, extended-value KL divergence~\eqref{eq:intro-kl} is jointly convex, whence convexity of $\Env_\eta[g]$ follows by the partial minimization rule. 
For~\eqref{eq:right-bm-prox-identity}, apply Danskin's theorem and the first-order optimality condition to~\eqref{eq:envelope} with~$x \in \ri(\Simp)$.
Finally, invoking~\eqref{eq:envelope} at~$x$ and~$y$, with~$u := \prox_{\eta g}(x)$ as a feasible point in the latter case, gives
%This gives
% as a feasible point in and subtract the identity defining $\Env_\eta[g](x)$ we get
\begin{equation}
\begin{aligned}
\label{eq:right-bm-one-step}
    \eta \Env_\eta[g](y)- \eta \Env_\eta[g](x)
    \,\le\,
    \Div(y,u)-\Div(x,u)
    \,=\,
    \Div(y,x)+ \left\langle \log x-\log u,y-x\right\rangle.
\end{aligned}
\end{equation}
We used the three-point identity in the final step.
Combining this with~\eqref{eq:right-bm-prox-identity} and~\eqref{eq:grad-tangent} gives~\eqref{eq:right-bm-relative-smoothness}.
%\begin{equation}
%\label{eq:right-bm-three-point}
%\end{equation}
\end{proof}

\iffalse
The scaling identity
\begin{equation}
\label{eq:right-bm-scaling}
    \Env_\eta[a g]
    =
    a\,\Env_{a\eta}[g],
    \qquad a>0,
\end{equation}
will be used to keep the lower-bound normalization separate from the final smoothness constant.  In particular,
\begin{equation}
\label{eq:right-bm-normalization}
    \frac12\,\Env_{1/(2L)}[L g]
    =
    \frac L2\,\Env_{1/2}[g].
\end{equation}
By Lemma~\ref{lem:envelope-regularity}, the left-hand side belongs to $\mathcal{F}_h(L)$.
\fi

Next, we establish {\em locality} of the smoothing mechanism based upon the right Bregman--Moreau envelope, ensuring that the envelopes of locally coinciding functions (locally) agree to first order.
%then their smoothed counterparts agree to first order.
%the locality property of the  

\begin{lemma}%[Locality of the right envelope]
\label{lem:right-bm-locality}
Let~$g,\tilde g:\Simp\to\mathbb{R}$ be finite, continuous, convex, and such that~$\tilde g\ge g$ on~$\Simp$. 
Fix arbitrary~$\hat x\in\ri(\Simp)$ and let~$\hat u:=\prox_{\eta g}(\hat x)$. 
If $\tilde g=g$ in a relative neighborhood of $\hat u$, then %one has
\begin{equation}
\label{eq:right-bm-locality-equality}
    \Env_\eta[\tilde g](x)=\Env_\eta[g](x)
    \qquad \forall x \in \mathcal{X}
\end{equation}
in some relative neighborhood $\mathcal{X}$ of $\hat x$. 
In particular,~$\bnabla\Env_\eta[\tilde g](\hat x)= \bnabla \Env_\eta[g](\hat x)$.
%In particular, the two envelopes have the same value and affine gradient at $\bar x$.
\end{lemma}

\begin{proof}
Let~$\mathcal{U}$ be a relative neighborhood of~$\hat u$ where~$\tilde g=g$. 
By continuity of~$\prox_{\eta g}(\cdot)$, see Lemma~\ref{lem:prox-mapping}, in some relative neighborhood $\mathcal{X}$ of $\hat x$ one has
\[
\prox_{\eta g}(x) \in \mathcal{U} \qquad \forall x \in \mathcal{X}.
\]  
Now, fix arbitrary~$x \in \mathcal{X}$ and put~$u := \prox_{\eta g}(x)$. 
Since $\tilde g \ge g$ everywhere on~$\Simp$, by~\eqref{eq:envelope} we get
\begin{equation*}
    \Env_\eta[\tilde g](x)\ge \Env_\eta[g](x).
\end{equation*}
On the other hand, by invoking~\eqref{eq:envelope} for~$\tilde g$ with~$u$ as a feasible point (note that~$u \in \mathcal{U}$), we get
% in the minization problem characterizing~$\Env_\eta[\tilde g](x)$, we get
\begin{equation*}
\begin{aligned}
    \Env_\eta[\tilde g](x)
    &\le
    \tilde g(u)+\eta^{-1}\Div(x,u)
    \\
    &=
    g(u)+\eta^{-1}\Div(x,u)
    =
    \Env_\eta[g](x).
\end{aligned}
\end{equation*}
%where the equality uses that~$u \in \mathcal{U}$. 
Thus~$\Env_{\eta} [\tilde g] = \Env_{\eta} [g]$ on~$\mathcal{X}$, and~$\bnabla\Env_\eta[\tilde g](\hat x)= \bnabla \Env_\eta[g](\hat x)$ follows from differentiability on $\ri(\Simp)$.
\end{proof}

%\input{old}
\subsection{Oracle complexity model and formal statement of the result}
\label{sec:formal}

We work in the oracle complexity model of~\cite{nemirovski1983problem}, adapted to account for the lack of differentiability on the boundary for objectives in~$\Fh(L)$. 
%Definition~\ref{def:entropy-smoothness}.
Let us give a brief yet rigorous summary of this model.

%In this model, 
Any~$f \in \Fh(L)$ specifies a {\em first-order oracle}~$\Orc_f$ which, when queried at any~$x \in \ri(\Simp)$, returns 
%the datum
\begin{equation}\label{eq:oracle}
\Orc_f(x) := (f(x), \bnabla f(x))
\end{equation}
where~$\bnabla f$ is the tangent gradient of~$f$. 
%in~\eqref{eq:grad-tangent}).
As we explained in the beginning of Section~\ref{sec:prelim}, the use of~$\bnabla f$ is without loss of generality, since~$f$ is only defined over~$\Simp$ whose affine hull is parallel to~$\Tan$. 
However, one could still ask if anything could be gained by allowing the objective to be defined on the whole orthant~$\R^d_{+}$, with access to the full gradient oracle, while only requiring the relative smoothness condition in~\eqref{eq:intro-sandwich} to hold on~$\Simp$.
To address this, in Appendix~\ref{app:conic} we show that~\eqref{eq:intro-lower} remains valid for minimization in the class of~$1$-homogeneous extensions~\cite{ovcharov2018proper} of functions in~$ \Fh(L)$.
%The underlying reason is that such objectives can be thought of as extensions of~$f \in \Fh$; as such, their gradient is defined by 
%some discussion.
%Its use merits a remark.

%\begin{remark}
%\label{rem:emulation}
%Using~$\bnabla f$ is natural, since~$f$ is defined on a set~$\Simp$ in the affine subspace parallel to~$\Tan$.
%\end{remark}


%Note that the oracle~$x \mapsto (f(x),\bnabla f(x))$ is {more powerful} than the oracle~$x \mapsto \nabla f^+(x)$.

A deterministic {\em $T$-step first-order method} (FOM) makes~$T$ sequential queries~$x_1, \dots, x_{T} \in \ri(\Simp)$ of such an~$\Orc_f$, and outputs~$x_{T+1} \in \Simp$. 
Thus, a~$T$-step FOM is specified by a collection of mappings
%comprised of (i) a rule for selecting the initial query~$x_0 \in \ri(\Simp)$; (ii) a collection of mappings
\begin{equation}
\label{eq:method-maps}
\begin{aligned}
\Phi_t: \Omega^{t-1}  &\to \ri(\Simp), \qquad t \in [T]; \\
\bar \Phi_{T+1}: \Omega^T &\to \Simp,
\end{aligned}
\end{equation}
where~$\Omega = \R \times \Tan$, and~$\Omega^0$ is a singleton (so~$\Phi_1$ merely selects a specific point~$x_1 \in \ri(\Simp)$).
We let~$\FOM_d(T)$ be the class of all~$T$-step FOMs, as per~\eqref{eq:method-maps}. 
When a given method~$\Mtd \in \FOM_d(T)$ is instantiated with an oracle~$\Orc_f$ associated to a specific instance~$f \in \Ent(L)$, the resulting sequence
%utilized to generate the subsequent queries~$x_{t} = \Phi_{t}(\Orc_f(x_0),\dots,\Orc_f(x_{t}))$ in iterations~$t = 1, \dots, T-1;$
%(iii) a mapping~$\Omega_T: \R^{T(d+1)} \to \Simp$ used to form the final output~$x_T$.
%When a given method~$\Mtd$ 
%\[
%x_0^{\vphantom\Mtd} = x_0^{\Mtd}, \;\; x_1^{\vphantom\Mtd} = x_1^{\Mtd}(f), \;\; \dots, \;\; x_T^{\vphantom\Mtd} = x_T^{\Mtd}(f),
%\]
%\[
%x_0^{\Mtd}; \;\; \Orc_f(x_0^{\Mtd}); \;\;  x_1^{\Mtd}(f) = \Phi_1(\Orc_f(x_0^{\Mtd})); \;\; \Orc_f(x_1^{\Mtd}(f)); \;\; \ldots \;\; x_{T}^{\Mtd}(f) = \bar\Phi_{T}(\Orc_f(x_0^\Mtd),\dots,\Orc_f(x_{T-1}^\Mtd))
%\]
\begin{equation}
\label{eq:method-iterates}
x_1^{\Mtd} = \Phi_1^{\vphantom\Mtd}, 
\;\;
x_2^{\Mtd}(f) = \Phi_2^{\vphantom\Mtd}(\Orc_f(x_1^{\Mtd})), 
%\quad x_2^{\Mtd}(f) = \Phi_2^{\vphantom\Mtd}(\Orc_f(x_0^{\Mtd}),\Orc_f(x_1^\Mtd(f))), \\
\;\; \ldots, \;\; 
x_{T+1}^{\Mtd}(f) = \bar\Phi_{T+1}^{\vphantom\Mtd}(\,\Orc_f(x_1^\Mtd),\; \dots,\; \Orc_f(x_{T}^\Mtd(f))),
\end{equation}
along with the responses~$\Orc_f(x_1^\Mtd),\Orc_f(x_2^\Mtd(f)),\dots,\Orc_f(x_{T}^\Mtd(f))$, is called {\em the transcript} of~$\Mtd$ on~$f$.
Note that~$\Mtd$ selects each~$x_t$ {before} receiving the oracle response~$\Orc_f(x_t)$, and the crucial property of a transcript is its {\em  consistency}: the oracle responses in~\eqref{eq:method-iterates} correspond to {the same}~$f \in \Fh(L)$. 


To quantify the hardness of first-order optimization over~$\Fh(L)$, we define its {\em minimax~$T$-risk}
%, this can be formalized as bounding from below {\em the minimax~$T$-risk of~$\Ent_d(L)$}, defined as
\begin{equation}
\label{eq:minimax-risk}
\Risk_{d}(T,L) \; 
:=
\inf_{\Mtd^{\vphantom{2^L}} \,\in\, \FOM_{d}(T)} 
\quad
\sup_{f \,\in\, \Fh(L)} 
\quad
\left\{ f(x_{T+1}^\Mtd(f))- \min_{x \in \Simp} f(x) \right\}.
\end{equation}
With this definition at hand, we are now in the position to give the rigorous statement of Theorem~\ref{th:intro-lower}.

\begin{theorem}
\label{th:lower}
For all~$L > 0$,~$T \ge 1$ and~$d \ge 8(T+1)^2 + T$, the minimax~$T$-risk of~$\Fh(L)$ satisfies
%deterministic local first-order method that
%For any deterministic first-order method that makes~$T-1$ queries~$x_{0}, \dots, x_{T-1} \in \ri(\Simp)$ to the oracle~$x \mapsto (f(x),\nabla f(x))$ and returns~$x_{T} \in \Simp$, there exists~$f \in \Fh(L)$ such that
\begin{equation}
\label{eq:lower}
\Risk_{d}(T,L) > \frac{L}{4(T+1)}. 
\end{equation}
\end{theorem}
%
To prove Theorem~\ref{th:lower}, we shall proceed via the {``resisting oracle''} approach: 
interacting with arbitrary~$\Mtd \in \FOM_d(T)$, we shall construct a sequence~$\omega_1, \dots, \omega_T \in \Omega$, with
$
\omega_{t} = \omega_t(x_1, \dots, x_t),
$
which, along with the associated sequence of queries
$
\smash{
x_1 = \Phi_1,\;
x_2 = \Phi_2(\omega_1),
\ldots,
x_{T+1} = \bar\Phi_{T+1}(\omega_1,\dots,\omega_{T}), 
}
$
corresponds to a consistent transcript, i.e.,~matches some~$f \in \Fh(L)$ in the sense that~$\omega_t = \Orc_f(x_t^\Mtd)$. 
%cf.~\eqref{eq:method-iterates}.
%The ``not-so-classical'' elements of the construction are the specific choice of the resisting oracle, and the smoothing method used.
%Showing that no~$T$-step FOM perform on~$\Fh(L)$ better than at a certain level reduces to constructing a suitable resisting oracle. 
%In what follows next, we describe the 
%Our next goal is to construct such a resisting oracle, 
Our hard instance will be a maximum of affine functions, smoothed via the right Bregman-Moreau envelope to put it in~$\Fh(L)$.
In what follows, we first describe the general family of such nonsmooth functions and study their envelopes, then present the adaptive construction and complete the proof.
%Before we present this, let us explain how such 

%reduces to constructing a resisting oracle such that

%{\em oracle-complexity lower bound}
%To show an oracle complexity lower bound, one must construct {a ``resisting oracle,''} meaning that one must describe how to select the next data 
 
%completed with a rule for selecting the initial~$x_0 \in \ri(\Simp)$. 
%Note that algorithms that only use~$\nabla f(x)$, but not~$f(x)$, are formally allowed in this model, and thus are subject to the lower bound we are about to present; similarly, the lower bound remains valid if an oracle only returns the tangent gradient~$\bnabla f$. 

%Let $\mathcal O$ be a local oracle on $\Delta_d$ with values in an arbitrary answer space $\mathfrak A$. Locality means that, whenever $f_1$ and $f_2$ coincide on a relative neighborhood of $x$ in $\Delta_d$, there holds
%\begin{equation}\label{eq:local-oracle-definition}
%\mathcal O(f_1,x)=\mathcal O(f_2,x).
%\end{equation}
%We use the following fixed-budget deterministic convention. A $T$-step method is a collection of maps
%\begin{equation}\label{eq:method-maps}
%\Phi_1\in \Delta_d,
%\qquad
%\Phi_t:\mathfrak A^{t-1}\to \Delta_d,\quad t=2,\dots,T.
%\end{equation}
%Given an objective $f$, the generated points are
%\begin{equation}\label{eq:method-recursion}
%x_1^M(f):=\Phi_1,
%\qquad
%x_t^M(f):=
%\Phi_t\bigl(
%\mathcal O(f,x_1^M(f)),\dots,\mathcal O(f,x_{t-1}^M(f))
%\bigr),
%\quad t=2,\dots,T,
%\end{equation}
%and the reported point is $x_T^M(f)$. Thus $x_T^M(f)$ is chosen before the oracle answer at $x_T^M(f)$ can be used. This is the indexing convention used throughout the adversarial construction.
%
%Let $\FOM_{d}(T)$ denote the class of deterministic~$T$-step FOMs satisfying \eqref{eq:method-maps}--\eqref{eq:method-recursion}. 
%{\em The minimax~$T$-risk of~$\Fh(L)$} is defined as
%\begin{equation}\label{eq:intro-risk}
%\Risk_{\mathcal F_h(L),\Delta_d,\mathcal O}(T):=
%\inf_{M\in \mathfrak M_T^{\det}}
%\sup_{f\in \mathcal F_h(L)}
%\bigl[f(x_T^M(f))-\Opt(f)\bigr].
%\end{equation}
%

Before we proceed, a simple remark is in order. The minimax~$T$-risk is~$1$-homogeneous in~$L$, i.e.
\begin{equation}
\label{eq:risk-homogeneity}
\Risk_d(T,L) = L\Risk_d(T,1),
\end{equation}
as seen by noting that~$f \in \Fh(1)$ implies~$Lf \in \Fh(L)$, and that passing from~$f$ to~$Lf$ does not influence consistency of a transcript.
Since the right-hand side of~\eqref{eq:lower} has the same homogeneity, it suffices to treat the~$L = 1$ case; in other words, exhibit~$f \in \Ent(1)$ that certifies~\eqref{eq:lower} with~$L = 1$.

%The hard instances are smoothings of scaled nonsmooth backbones, the smoothing being carried out by the right Bregman--Moreau envelope of Section~\ref{sec:envelope}. Set
%\begin{equation}
%        \theta:=\tfrac12,
%        \qquad
%        \eta_L:=\tfrac1{2L},
%\end{equation}
%so that $\theta$ is the dimensionless envelope scale and $\eta_L$ the scale producing the target smoothness constant. For every convex, continuous $g:\Simp\to\R$, the scaling identity \eqref{eq:right-bm-scaling} specializes to the normalization \eqref{eq:right-bm-normalization}, which in the present notation reads
%\begin{equation}\label{eq:scaling-id}
%        \tfrac12\,\Env_{\eta_L}[Lg]
%        =
%        \tfrac L2\,\Env_\theta[g].
%\end{equation}
%Recall from Lemma~\ref{lem:envelope-regularity} that $\Env_{\eta_L}[Lg]$ is convex and $\eta_L^{-1}=2L$-smooth relative to $h$, in the sense of \eqref{eq:right-bm-relative-smoothness}; together with \eqref{eq:right-bm-normalization} this places the left-hand side of \eqref{eq:scaling-id}, and hence the function $(L/2)\Env_\theta[g]$ that we shall actually analyze, in the class $\Fh(L)$. We may therefore work throughout with the dimensionless envelope $\Env_\theta$ and reinstate the factor $L/2$ only at the very end.


\subsection{Hard instances and their properties}
\label{sec:hard-subclass}
%\section{Negative-coordinate backbones}

%\odima{TODO: we can use~$j$ instead of~$j_k$: the exposed coordinates do not repeat (correct in the next pass)}\\
Our hard instances are based on nonsmooth functions that are maxima of shifted coordinate atoms
%
%All backbones below are maxima of negative-coordinate atoms,
\begin{equation}
\label{eq:backbone}
        g(u)=\max_{k \in [K]}\bigl\{-(u)_{j_k} - b_k\bigr\},
\end{equation}
indexed by a finite set~$[K]$, with distinct coordinates $j_k$ and offsets~$b_k \ge 0$. 
For brevity, we shall refer to such functions as {\em backbones}.
In a backbone, each atom is~$1$-Lipschitz in~$\ell_\infty$-norm, and so is the whole backbone; this will be used later on.
By definition, the set of {\em $g$-active atoms} at~$u \in \ri(\Simp)$ is
\[
A_g(u):=\bigl\{k\in [K]:\ -(u)_{j_k}-b_k =  g(u)\bigr\},
\]
and~$\smash{\{j_k: k \in A_g(u)\}}$ are {\em $g$-active coordinates~at~$u$}.
%the set of atoms active at $u$, and call $i_r$, $r\in I(u)$, the active coordinates at $u$. 
To obtain a legitimate~$\smash{f \in \Fh(1)}$, we take the right Bregman--Moreau envelope of a backbone~$g$.
Indeed, by Lemma~\ref{lem:envelope-regularity} we have~$\smash{\eta S_\eta[g](\cdot) \in \Fh(1)}$ for any~$\eta > 0$; later on, we shall fix~$\eta = {1}/{2}$.

%Recall that the right Bregman--Moreau envelope of $g$ at scale $\theta$ is $\Env_\theta[g](x)=\min_{u\in\Simp}\{g(u)+\theta^{-1}\Div(x,u)\}$, cf.~\eqref{eq:right-bm-relative-smoothness}. 
%For $x\in\Simp$, the minimum is attained at a unique point of the relative interior of the face $\Delta_{\supp x}$, since $\Div(x,\cdot)$ is strictly convex there; we denote this minimizer by
%\begin{equation}
%\label{eq:prox-def}
%        u_g(x):=\argmin_{u\in\Simp}\bigl\{g(u)+\theta^{-1}\Div(x,u)\bigr\},
%\end{equation}
%
%so that $\Env_\theta[g](x)=g(u_g(x))+\theta^{-1}\Div(x,u_g(x))$. 
%On $\ri(\Simp)$ this is the prox-mapping $\prox_{\theta g}$ of Lemma~\ref{lem:prox-mapping}, which is continuous. 
%
Next, we prove some regularity properties of the right Bregman-Moreau envelope of a backbone.
We begin with a lemma that controls the entrywise growth of the prox-mapping associated with~$\eta g$.
%(ii) a lower bound for~$\smash{S_\eta[g]}$ in terms of~$g$.

\begin{lemma}
%[Entrywise growth]
\label{lem:entry-growth}
Let $g$ be as in~\eqref{eq:backbone}, and~$\eta < 1$. 
For all~$x\in\ri(\Simp)$ and~$j \in [d]$, 
%one has
%the prox-mapping~$\prox_{\eta g}(x)$ satisfies
%Then $\supp u=\supp x$, and
\begin{equation}
\label{eq:entry-growth}
0 < \eta (\prox_{\eta g}(x))_j 
\le \frac{\eta}{1-\eta} (x)_j.
\end{equation}
In particular, for~$\eta \le \frac{1}{2}$ one has~$0 < \eta \prox_{\eta g}(x) \le x$ entrywise.
\end{lemma}

%\odima{HERE}

\begin{proof}
Let~$\hat u := \prox_{\eta g}(x)$.
The left inequality in~\eqref{eq:entry-growth}, i.e.~that~$\prox_{\eta g}(x) \in \ri(\Simp)$, is immediate: 
%from~\eqref{eq:envelope}--\eqref{eq:prox} and~\eqref{eq:intro-kl}: 
if~$\hat u$ has a zero entry, then~$\Div(x,\hat u) = +\infty$ by~\eqref{eq:intro-kl}, which contradicts the optimality of~$\hat u$. 
%
For the right inequality, since the minimum in~\eqref{eq:envelope} is attained on~$\smash{\ri(\Simp)}$, the optimality condition writes as
\begin{equation}
\label{eq:optimality-for-g}
\exists \mu \in \R: \quad \partial g(\hat u) \ni \eta^{-1} \frac{x}{\hat u} - \mu \ones,
\end{equation}
where the division is entrywise.
Now, let~$(j_k)_{k \in A_g(\hat u)}$ be the active part of the finite sequence~$(j_k)_{k \in [K]}$ in~\eqref{eq:backbone}, 
and consider arbitrary subgradient
$\xi(\hat u) = -\sum_{k \in A_g(\hat u)} w_k e_{j_k}$ of~$g(\cdot)$ at~$\hat u$, where~$w_k \ge 0$ and~$\smash{\sum_{k \in A_g(\hat u)} w_k = 1}$. 
Letting~$\lambda_i$, for~$i \in [d]$, be the sum of all active multipliers for each coordinate~$i$,
\[
\lambda_i := \sum_{k \,\in\, A_g(\hat u)} w_k\, \delta_{i j_k},
%\quad i \in [d].
\]
we can write~$\xi(\hat u) = -\sum_{i \in [d]} \lambda_i e_i$ where~$(\lambda_1, \dots, \lambda_d) \in \Simp$ (in particular,~$\lambda_i \le 1$).
Plugging in~\eqref{eq:optimality-for-g} the relevant subgradient~$\xi(\hat u)$---i.e., one realizing the inclusion in~\eqref{eq:optimality-for-g}---and rearranging, we get
\[
\frac{(x)_i}{(\hat u)_i} = \eta (\mu -\lambda_i), \quad i \in [d].
\]
%where we write~$x_i = \langle x, e_i \rangle$ for brevity (and similarly for~$\hat u_i$).
Multiplying by~$(\hat u)_i$, summing over~$i$, and using that~$x,\hat u \in \Simp$, we see that~$\eta(\mu - \hat m_{\lambda}) = 1$, where~$\hat m_{\lambda} := \sum_{i \in [d]} \lambda_i (\hat u)_i$ 
is the $\lambda$-weighted average of the entries of~$\hat u$.
Combining with the above, 
%Combining with the above expression, we get
\begin{equation}
\label{eq:ratio}
\frac{(x)_i}{(\hat u)_i} = 1 + \eta(\hat m_{\lambda} - \lambda_i), \quad i \in [d].
\end{equation}
Noting that~$\hat m_{\lambda} \ge 0$ and~$\lambda_i \le 1$, we get~$(\hat u)_i \le (1-\eta)^{-1} (x)_i$, that is the right inequality in~\eqref{eq:entry-growth}.
\end{proof}

%%%%%%%%%%%%%%%%%%%%%%%% Active face proof %%%%%%%%%%%%%%%%%%%%%%%%

\begin{comment}
\proofstep{1}.
Set $J:=\supp x$. Since $\Env_\theta[g](x)$ is finite, attained for instance at $u=x$, the minimizer satisfies $\Div(x,u)<\infty$, so $J\subseteq\supp u$ by \eqref{eq:intro-kl}. Suppose, for contradiction, that $u$ places mass $m:=\sum_{j\notin J}u_j>0$ outside $J$. Renormalize onto $J$ by setting $\bar u_j:=u_j/(1-m)$ for $j\in J$ and $\bar u_j:=0$ otherwise; then $\bar u\in\Simp$, and $\sum_{j\in J}u_j=1-m$ gives $\norm{u-\bar u}_\infty\le m$. By $1$-Lipschitzness of $g$,
\begin{equation}\label{eq:renorm-lip}
        g(\bar u)\le g(u)+m.
\end{equation}
Moreover, because $x$ is supported on $J$ and $\sum_{j\in J}x_j=1$,
\begin{equation}\label{eq:renorm-kl}
        \Div(x,u)=\Div(x,\bar u)-\log(1-m).
\end{equation}
Combining \eqref{eq:renorm-lip} and \eqref{eq:renorm-kl}, the change in the objective of \eqref{eq:prox-def} satisfies
\begin{align}
\bigl[g(\bar u)+\theta^{-1}\Div(x,\bar u)\bigr]
-\bigl[g(u)+\theta^{-1}\Div(x,u)\bigr]
&\le
m+\theta^{-1}\log(1-m)\notag\\
&<
m\bigl(1-\theta^{-1}\bigr)
<0,
\end{align}
where the strict inequality uses $\log(1-m)<-m$, and the last step uses $\theta=\tfrac12<1$. This contradicts the optimality of $u$. Hence $\supp u=J$, and in particular $u\in\ri(\Delta_J)$.

\proofstep{2}.
It remains to bound $u_j$ for $j\in J$. The first-order optimality condition for \eqref{eq:envelope} on the face $\Delta_J$ reads $0\in\partial g(u)+\eta^{-1}\nabla_u\Div(x,u)+\spanof{\ones}$. Choose a subgradient $v=-\sum_{r\in I(u)}\lambda_r e_{i_r}\in\partial g(u)$ with $\lambda_r\ge0$ and $\sum_{r\in I(u)}\lambda_r=1$, and write $\lambda_j:=\sum_{r\in I(u):\,i_r=j}\lambda_r$ for the active multiplier carried by coordinate $j\in J$; active atoms whose coordinate lies outside $J$ are constant on $\Delta_J$ and drop out. Since $\nabla_{u_j}\Div(x,u)=-x_j/u_j$, the optimality condition is, with multiplier $\mu\in\R$ for the constraint $\sum_{j\in J}u_j=1$,
\begin{equation}
        -\lambda_j-\eta^{-1}\frac{x_j}{u_j}+\mu=0,
        \qquad j\in J.
\end{equation}
Equivalently $x_j/u_j=\eta(\mu-\lambda_j)$. Multiplying by $u_j$, summing over $j\in J$, and using $\sum_{j\in J}x_j=\sum_{j\in J}u_j=1$, we obtain $1=\eta(\mu-m_\lambda)$ with $m_\lambda:=\sum_{j\in J}\lambda_j u_j$, so that $\mu=\eta^{-1}+m_\lambda$ and
\begin{equation}\label{eq:ratio}
        \frac{x_j}{u_j}=1+\eta(m_\lambda-\lambda_j),
        \qquad j\in J.
\end{equation}
Since $m_\lambda\ge0$ and $\lambda_j\le1$, the right-hand side of \eqref{eq:ratio} is at least $1-\eta$. Therefore $\eta u_j\le\tfrac{\eta}{1-\eta}x_j$ for $j\in J$, while for $j\notin J$ both sides of \eqref{eq:entry-growth} vanish; this proves \eqref{eq:entry-growth}, and the case $\eta\le\tfrac12$ follows since then $\tfrac{\eta}{1-\eta}\le1$.
\end{comment}

%%%%%%%%%%%%%%%%%%%%%%%%%%%%%%%%%%%%%%%%%%%%%%%%%%%%%%%%%%%%%%%%%%%


The next lemma controls the value decrease of a backbone caused by smoothing.
% the B  by the largest active coordinate.


\begin{lemma}
%[Smoothing gap]
\label{lem:loss}
Let~$g$ be as in~\eqref{eq:backbone},~$x\in\ri(\Simp)$, % $x \in \Simp$ in the active-face formulation
and~$\hat u =\prox_{\eta g}(x)$. 
Furthermore, suppose that~$(\hat u)_{j_k} \le M$ for all~$g$-active coordinates at~$\hat u$, i.e.~for all~$\smash{k \in A_g(\hat u)}$.
Then
\begin{equation}
\label{eq:loss-bound}
        \Env_\eta[g](x)\ge g(x)-\eta M.
\end{equation}
\end{lemma}

\begin{proof}
Rearranging \eqref{eq:ratio} gives $(x)_i-(\hat u)_i=\eta (\hat m_\lambda-\lambda_i) (\hat u)_i$ in terms of~$\hat m_{\lambda} = \sum_{i \in [d]} \lambda_i (\hat u)_i$, the quantity introduced in the proof of Lemma~\ref{lem:entry-growth}. 
Whence
\begin{equation}
\label{eq:displacement}
        \abs{(x)_i-(\hat u)_i}=\eta \abs{\hat m_\lambda-\lambda_i} (\hat u)_i \, .
\end{equation}
Since~$\hat m_\lambda$ is a convex combination of the active-coordinate values~$\{(\hat u)_{j_k}, k \in A_g(\hat u)\}$, we have~$\hat m_\lambda\le M$ by the premise of the lemma.
Let us show that
\begin{equation}
\label{eq:uniform-bound}
\| x - \hat u \|_{\infty} \le \eta M.
\end{equation}
To that end, since~$\hat m_\lambda, \lambda_i \in [0,1]$,
we get~$\abs{(x)_i - (\hat u)_i} \le \eta M$ for active coordinates~$i \in \{j_k, k \in A_g(\hat u)\}$. %cf.~\eqref{eq:displacement}.
Meanwhile, for inactive coordinates~$\lambda_i = 0$, so~\eqref{eq:displacement} and~$(\hat u)_i \le 1$ result in
$
\abs{(x)_i - (\hat u)_i} \le \eta \hat m_{\lambda} \le \eta M.
$
This verifies~\eqref{eq:uniform-bound}. 
Now, by the $1$-Lipschitzness of a backbone, it follows that
$
g(\hat u)\ge g(x)-\eta M.
$
In turn, this implies~$\Env_\eta[g](x)=g(\hat u)+\eta^{-1}\Div(x,\hat u) \ge g(\hat u) \ge g(x)-\eta M$, as claimed in~\eqref{eq:loss-bound}.
\end{proof}

%%%%%%%%%%%%%%%%%%%%%%%% Active face proof %%%%%%%%%%%%%%%%%%%%%%%%

\begin{comment}
Rearranging \eqref{eq:ratio} gives $x_j-u_j=\eta u_j(m_\lambda-\lambda_j)$, hence
\begin{equation}\label{eq:displacement}
        \abs{x_j-u_j}=\eta\, u_j\,\abs{m_\lambda-\lambda_j},
        \qquad j\in J:=\supp x.
\end{equation}
Note that $m_\lambda=\sum_{r\in I(u)}\lambda_r u_{i_r}$ is a convex combination of active-coordinate values, so $m_\lambda\le M$; moreover $0\le m_\lambda,\lambda_j\le1$. If $\lambda_j>0$, then $j$ is an active coordinate, so $u_j\le M$, and \eqref{eq:displacement} gives $\abs{x_j-u_j}\le\eta u_j\le\eta M$ since $\abs{m_\lambda-\lambda_j}\le1$. If $\lambda_j=0$, then $\abs{x_j-u_j}=\eta u_j m_\lambda\le\eta M$, using $u_j\le1$ and $m_\lambda\le M$. In either case $\abs{x_j-u_j}\le\eta M$, and the same bound holds trivially for $j\notin J$. Hence $\norm{x-u}_\infty\le\eta M$, and by $1$-Lipschitzness of $g$ it follows that $g(u)\ge g(x)-\eta M$. Adding the nonnegative term $\eta^{-1}\Div(x,u)$ yields $\Env_\eta[g](x)=g(u)+\eta^{-1}\Div(x,u)\ge g(x)-\eta M$, which is \eqref{eq:loss-bound}.
\end{comment}

%%%%%%%%%%%%%%%%%%%%%%%%%%%%%%%%%%%%%%%%%%%%%%%%%%%%%%%%%%%%%%%%%%%
\section{Proof of Theorem~\ref{th:intro-lower}}
\label{sec:proof}
%
In this section, we prove Theorem~\ref{th:lower}, and Theorem~\ref{th:intro-lower} along with it.
To that end, we first implement the resisting oracle announced in Section~\ref{sec:formal}: interacting with an arbitrary method~$\Mtd\in\FOM_d(T)$, cf.~\eqref{eq:method-maps}--\eqref{eq:method-iterates}, after each query we add a new atom at the currently smallest-mass coordinate, with offset increased in constant increments. 
Smoothing via the right Bregman--Moreau envelope with~$\eta = \half$, while producing an instance in~$\Fh(1)$ by Lemma~\ref{lem:envelope-regularity}, ensures that a new atom is strictly dominated near all the previous queries; this effect is attained through~Lemma~\ref{lem:loss} and offsets. 
This guarantees transcript consistency: the final instance matches the earlier oracle answers.
%the finally exhibited instance matches the earlier oracle answers.

%this ensures consistency of the transcript.
%As the result, the oracle cannot distinguish the running objective from the single instance assembled at the end, and consistency of the transcript follows. 
%Recall that it suffices to certify~\eqref{eq:lower} for~$L=1$; we accordingly fix the envelope scale~$\eta:=\half$ promised in Section~\ref{sec:hard-subclass}, so that~$\eta\Env_\eta[g]\in\Fh(1)$ for every backbone~$g$ by Lemma~\ref{lem:envelope-regularity}. 

The argument proceeds in three stages. %presented in subsequent subsections.
In Section~\ref{sec:adaptive-backbone}, we detail the construction. 
In Section~\ref{sec:transcript-consistency}, we prove a locality lemma ensuring that new atoms are dominated by the previous ones, and deduce transcript consistency. 
% derive transcript consistency as a result.
In Section~\ref{sec:gap-proof}, we bound the suboptimality gap; this is done by augmenting the exposed face with an extra dimension and using the center of the augmented face.
%that lives in a ``hidden'' subspace.
% consistency proposition that proves  
%closing the resisting-oracle loop (Proposition~\ref{prop:transcript}); the optimality gap is then estimated in Section~\ref{sec:gap-proof}.

%\odima{HERE}

\subsection{Constructing the hard instance}
\label{sec:adaptive-backbone}


Recall the assumption~$d \ge 8(T+1)^2+T$ in the premise. Throughout, fix~$\Mtd\in\FOM_d(T)$ and set
\begin{equation}
\label{eq:params}
        \veps_T :=\frac{1}{8(T+1)^2}\,,
        \qquad
        \rho:=\frac{2}{3}\,,
        \qquad
        \delta_T :=(2+\rho)\veps_T =\frac{1}{3(T+1)^2}\,.
        \qquad
\end{equation}
We shall select distinct coordinates~$j_1,\dots,j_{T},j_{T+1}$ incrementally, in response to~$\Mtd$'s queries. 

{\em Rounds~$t\in [T]$.} Once~$\Mtd$ has issued a query~$x_t\in\ri(\Simp)$ determined via~\eqref{eq:method-iterates} by the previous answers, 
we construct the answer to~$x_t$ by selecting a coordinate where~$x_t$ has the smallest mass:
\begin{equation}
\label{eq:fresh}
        j_t \in \Argmin_{i \,\in\, [d] \setminus E_{t}} \; \langle x_t, e_i \rangle, \qquad E_{t} := \{j_1,\ldots,j_{t-1}\},
\end{equation}
%on which~$x_t$ carries little mass. 
%
where~$E_{1} = \emptyset$ by convention.
We then let
\begin{equation}
\label{eq:atoms}
        a_t(x):=-\langle x, e_{j_t} \rangle -(t-1)\delta_T, \qquad 
        g_t(x) :=\max_{s \in [t]} a_s(x),
        \qquad
        f_t := \tfrac{1}{2}\Env_{1/2}[g_t]
\end{equation}
and return
%using~$f_t(\cdot)$ above, 
\begin{equation}
\label{eq:oracle-answers}
\Orc_{f_t}(x_t) = (f_t(x_t), \nabla f_t(x_t)),
\end{equation}
cf.~\eqref{eq:oracle}, as the answer to query~$x_t$. 
Note that~$g_t$, cf.~\eqref{eq:atoms}, is a backbone in the sense of~\eqref{eq:backbone}. 
This procedure specifies the resisting oracle at rounds~$t \in [T]$, corresponding to~$\Mtd$'s queries~$x_1, \dots, x_T$.

{\em Final round.}
To finalize the construction and exhibit the hard instance, we use the candidate minimizer~$x_{T+1} \in \Simp$ returned by~$\Mtd$ after the final oracle call.
Namely, we invoke~\eqref{eq:fresh} with~$t = T+1$, producing
\[
j_{T+1} \in \Argmin_{i \,\in\, [d] \setminus E_{T}} \; \langle x_{T+1}, e_i \rangle;
%\qquad E_{T} := \{j_1,\ldots,j_{T}\};
\]
we then let~$g_{T+1}(x) := \max\{g_T(x), a_{T+1}(x)\}$ with~$a_{T+1}(x):=-\langle x, e_{j_{T+1}} \rangle - T\delta_T$, cf.~\eqref{eq:atoms}, and take
\begin{equation}
\label{eq:final-instance}
f_{T+1} := \tfrac{1}{2} \Env_{1/2}[g_{T+1}]
\end{equation}
as the hard instance. Notice that~$g_{T+1}$ is still a backbone, therefore~$f_{T+1} \in \Fh(1)$ by Lemma~\ref{lem:envelope-regularity}.

Before we begin the proof of transcript consistency, let us record one implication of Lemma~\ref{lem:entry-growth}.

\begin{proposition}
\label{prop:fresh-entries}
For every~$t \in [T]$, selection rule~\eqref{eq:fresh} guarantees the following for~$\veps_T$ as in~\eqref{eq:params}:
\begin{equation}
\label{eq:fresh-entries}
\langle \prox_{\half g_{t}}(x_t), e_{j_t} \rangle 
\le 2\langle x_t, e_{j_t} \rangle 
\le 2\veps_T.
\end{equation}
Moreover, the right-hand inequality in~\eqref{eq:fresh-entries} extends to~$t = T+1$, i.e.~it holds that~$\langle x_{T+1}, e_{j_{T+1}} \rangle \le \veps_T.$ 
\end{proposition}
%
\begin{proof}
For any~$t \in [T+1]$, the cardinality of exposed support~$E_t = \{j_1,\dots,j_{t-1}\}$ is at most~$T$, so that of~$H_t := [d] \setminus E_t$ is at least
$
d-T \ge \veps_T{}^{-1},
$
cf.~\eqref{eq:params}. 
As such,~$\min_{i \in H_t} \langle x_t, e_i \rangle > \veps_T$ would imply
\[
\langle x_t, \ones \rangle \ge \sum_{i \in H_t} \langle x_t, e_i \rangle \ge |H_t| \, \min_{i \in H_t} \, \langle x_t, e_i \rangle > 1,
\]
contradicting~$x_t \in \Simp$. This proves the right-hand inequality in~\eqref{eq:fresh-entries} for all~$t \in [T+1]$. 
For~$t \in [T]$, we additionally have~$x_{t} \in \ri(\Simp)$, and the left-hand inequality in~\eqref{eq:fresh-entries} follows by Lemma~\ref{lem:entry-growth}.
%by~\eqref{eq:params}. 
%Were~$x_{t,j}>\veps$ for every~$j\in A_t$, summing over~$A_t$ would give~$\sum_{j\in A_t}x_{t,j}>\veps\,|A_t|\ge\veps\cdot\veps^{-1}=1$, contradicting~$\sum_{j=1}^d x_{t,j}=1$ with~$x_{t,j}\ge0$; hence some~$j\in A_t$ obeys~$x_{t,j}\le\veps$, and we put~$i_t:=j$.
\end{proof}
% The quadratic dimension~$d=\Omega(T^2)$ is exactly the slack that keeps this selection feasible through all~$T+1$ rounds. 
%With Proposition~\ref{prop:pigeonhole} at hand, 


% being the maximum of the~$t$ atoms~$\ell_1,\dots,\ell_t$ on the distinct coordinates~$i_1,\dots,i_t$ with nonnegative offsets~$(s-1)\delta$. Lemma~\ref{lem:entry-growth}, applied to~$g_t$ at~$x_t\in\ri(\Simp)$ with~$\eta=\half<1$, gives~$\eta\prox_{\eta g_t}(x_t)\le x_t$ entrywise, since~$\tfrac{\eta}{1-\eta}=1$ at~$\eta=\tfrac12$; equivalently~$\prox_{\eta g_t}(x_t)\le 2x_t$. Reading the~$i_t$-th coordinate and invoking the choice~\eqref{eq:fresh},
%\begin{equation}
%\label{eq:own-small}
%        u_{t,i_t}\le 2x_{t,i_t}\le 2\veps.
%\end{equation}
%After the answer at~$x_T$, the method reports a point~$x_{T+1}\in\Simp$ through~\eqref{eq:method-iterates}. Choose one further coordinate~$i_{T+1}\notin\{i_1,\dots,i_T\}$ with~$x_{T+1,i_{T+1}}\le\veps$, possible by the same counting since the forbidden set now has exactly~$T$ elements and~$d-T=\veps^{-1}$, and define~$\ell_{T+1}$ and~$g_{T+1}$ by~\eqref{eq:atoms}. The final hard instance is
%\begin{equation}
%\label{eq:final-instance}
%        f:=\eta\Env_\eta[g_{T+1}]=\tfrac12\,\Env_{1/2}[g_{T+1}],
%\end{equation}
%which lies in~$\Fh(1)$ by Lemma~\ref{lem:envelope-regularity}; cf.~the discussion in Section~\ref{sec:hard-subclass}.

\subsection{Ensuring transcript consistency}
\label{sec:transcript-consistency}

In this section, we establish transcript consistency for the construction presented in Section~\ref{sec:adaptive-backbone}.
The crux of the argument is that the ``fresh'' atom introduced at round~$t$ remains strictly dominated (``locally invisible'') in a vicinity~$\mathcal{U}_s$ of~$\smash{\hat u_s = \prox_{\half g_s}(x_s)}$; thus,~$g_t = g_s$ over~$\mathcal{U}_s$. 
By Lemma~\ref{lem:right-bm-locality}, this gives local coincidence of~$f_s = S_{1/2}[g_s]$ and~$f_{T+1} = S_{1/2}[g_{T+1}]$ in some relative neighborhood of~$x_s$.
%This ensures that  atoms do not perturb the backbone~$g_s$, and ultimately with~$g_{T+1}$. 
%\odima{The argumented is further split into two steps: a chronological margin at the prox points~$u_s$, transferred by Lipschitz continuity to a whole neighborhood, followed by the locality of the right envelope, Lemma~\ref{lem:right-bm-locality}.}

\begin{lemma}
%[``Local invisibility'']
\label{lem:local-invisibility}
For all~$1\le s<t\le T+1$, the following holds.
\vspace{-0.1cm}
\begin{enumerate}
\item 
The point~$\hat u_s = \prox_{\half g_s}(x_s)$ satisfies\vspace{-0.2cm}
\begin{equation}
\label{eq:margin}
         a_t(\hat u_s) \le g_s(\hat u_s) - \rho\veps_T. \vspace{-0.1cm}
\end{equation}
\item As a result, there is a relative neighborhood~$\mathcal{X}_s \subseteq\ri(\Simp)$ of~$x_s$ where one has~$f_t(x) = f_s(x)$.
\end{enumerate}
\end{lemma}

\begin{proof}
\proofstep{1}. %~{\em Chronological margin.} 
By~\eqref{eq:atoms}, a backbone dominates its last atom, so
$
g_s(\hat u_s) \ge a_s(\hat u_s) = -\langle \hat u_{s}, e_{j_s} \rangle - (s-1)\delta_T.
$
When combined with the bound~$\langle \hat u_{s}, e_{j_s} \rangle \le 2\veps_T$ furnished by Proposition~\ref{prop:fresh-entries}, cf.~\eqref{eq:fresh-entries}, this gives
\[
g_s(\hat u_s)\ge-2\veps_T-(s-1)\delta_T.
\] 
On the other hand, since~$\hat u_{s}$ has nonnegative entries,~$a_t(\hat u_s)=-\langle \hat u_{s}, e_{j_t}\rangle - (t-1)\delta_T \le-(t-1)\delta_T$. 
Subtracting the two estimates and using the facts that~$t > s$ and~$\delta_T=(2+\rho)\veps_T$, cf.~\eqref{eq:params}, we get
\begin{equation*}
        g_s(\hat u_s)-a_t(\hat u_s)
        \;\ge\;
        (t-s)\delta_T-2\veps_T
        \;\ge\;
        \delta_T-2\veps_T
%        \;=\; (2+\rho)\veps_T-2\veps_T
        \;=\; \rho\veps_T.
\end{equation*}
%that is~\eqref{eq:margin}. 

\noindent\proofstep{2}. %~{\em Local invisibility.} 
Our plan is to invoke Lemma~\ref{lem:right-bm-locality}.
Fix~$s<t$ and consider the neighborhood~$\mathcal{U}_s$ of~$\hat u_s$ as follows:
\[
\mathcal{U}_s 
:= 
\left\{
	u \in \ri(\Simp): \|u - \hat u_s\|_{\infty} \le \frac{\rho\veps_T}{4}
\right\}.
\]
Combining~\eqref{eq:margin} with~$1$-Lipschitzness of~$a_t(\cdot)$ and~$g_s(\cdot)$ in~$\ell_\infty$-norm, one has for all~$u \in \mathcal{U}_s$ and~$t > s$:
\[
        g_s(u)-a_t(u)
        \ge
        g_s(\hat u_s)-a_t(\hat u_s)-\frac{\rho\veps_T}{2}
        \stackrel{\eqref{eq:margin}}{\ge}
        \rho\veps_T - \frac{\rho\veps_T}{2}
        =\frac{\rho\veps_T}{2}
        >0.
\]
%Each atom~$a_t(\cdot)$ has constant gradient~$-e_{j_t}$, whose~$\ell_1$-norm---dual to~$\norm{\cdot}_\infty$ and so is~$1$-Lipschitz in~$\norm{\cdot}_1$. 
%being a pointwise maximum of such atoms, the backbone~$g_s=\max_{r\le s}\ell_r$ is~$1$-Lipschitz in~$\norm{\cdot}_\infty$ as well, cf.~the remark following~\eqref{eq:backbone}. Hence, for any~$w\in\Simp$ with~$\norm{w-u_s}_\infty\le\rho\veps/4$ and any~$t$ with~$s<t\le T+1$, the Lipschitz bounds~$\abs{g_s(w)-g_s(u_s)}\le\rho\veps/4$ and~$\abs{\ell_t(w)-\ell_t(u_s)}\le\rho\veps/4$ combine with the margin~\eqref{eq:margin} to yield
Thus, all future atoms~$a_t$,~$t > s$, are strictly dominated by~$g_s$ in the relative neighborhood~$\mathcal{U}_s$ of~$\hat u_s$ (note that we used that~$\hat u_s$ is in the relative interior of~$\Simp$). 
%the corresponding proximal point.
Applying this to~$t = s+1, \ldots, T+1$, 
\begin{equation*}
        g_{T+1}(u)
        =\max \left\{g_s(u), \max_{s < t \le T+1} a_t(u)\right\}
%        \Bigl(\max_{r\le s}\ell_r,\ \max_{s<t\le T+1}\ell_t\Bigr)
%        =\max\Bigl(g_s,\ \max_{s<t\le T+1}\ell_t\Bigr)
        =g_s(u) 
      \qquad \forall u \in \mathcal{U}_s.
\end{equation*}
Meanwhile, we have~$g_{T+1} \ge g_s$ pointwise on~$\Simp$, directly from the definitions of~$g_{T+1}$ and~$g_s$, cf.~\eqref{eq:atoms}.
%Since~$\rho\veps/4>0$ and~$u_s\in\ri(\Simp)$, the ball~$B$ is a genuine relative neighborhood of~$u_s=\prox_{\eta g_s}(x_s)$. 
%The pair~$g_s,g_{T+1}$ is finite, continuous and convex on~$\Simp$ (backbones, \eqref{eq:backbone}); one has~$g_{T+1}\ge g_s$ on all of~$\Simp$, being a maximum over a larger index set, with equality~$g_{T+1}=g_s$ on the relative neighborhood~$B$ of~$\prox_{\eta g_s}(x_s)$. 
As such, we may invoke Lemma~\ref{lem:right-bm-locality} with~$g=g_s$,~$\tilde g=g_{T+1}$,~$\hat x=x_s$ and~$\hat u = \hat u_s$; 
this gives a relative neighborhood~$\mathcal{X}_s \subseteq \ri(\Simp)$ in which~$S_{1/2}[g_{T+1}] = S_{1/2}[g_s]$, that is~$f_{T+1} = f_s$ (cf.~\eqref{eq:fresh}--\eqref{eq:final-instance}).
%furnishes a relative neighborhood~$V_s\subseteq\ri(\Simp)$ of~$x_s$ on which~$\Env_\eta[g_{T+1}]=\Env_\eta[g_s]$. Multiplying by~$\eta$ and recalling~\eqref{eq:atoms} and~\eqref{eq:final-instance} gives~$f=f_s$ on~$V_s$.
\end{proof}


\begin{proposition}[Transcript consistency]
\label{prop:transcript}
The queries~$x_1,\dots,x_T$, the oracle answers returned in~\eqref{eq:oracle-answers}, and the reported point~$x_{T+1}$ coincide with the transcript of~$\Mtd$ run on the objective~$f_{T+1}$.
\end{proposition}
%
\begin{proof}
Fix any~$s\in [T]$. 
By Lemma~\ref{lem:local-invisibility} with~$t = T+1$, we get~$f_{T+1}=f_s$ in a relative neighborhood~$\mathcal{X}_s$ of~$x_s$; in particular~$f_{T+1}(x_s)=f_s(x_s)$. 
Moreover, since the tangent gradient of~$f \in \Fh(1)$ at~$x_s$ is determined by the values of~$f$ in a~$\Tan$-neighborhood of~$x_s$ (cf.~\eqref{eq:grad-tangent}), we have~$\bnabla f_{T+1}(x_s)=\bnabla f_s(x_s)$ and
%For the differential datum, note first that the affine gradient~$\bnabla f(x_s)$ is determined by the values of~$f$ in a~$\Tan$-neighborhood of~$x_s$, cf.~\eqref{eq:grad-tangent}, whence~$\bnabla f(x_s)=\bnabla f_s(x_s)$; this is also the ``in particular'' clause of Lemma~\ref{lem:right-bm-locality}. To control the full gradient that the oracle~\eqref{eq:oracle} returns, recall from Section~\ref{sec:envelope} that the right envelope admits a canonical~$C^1$ extension to~$\R^d_{++}$, whose gradient at an interior point is an explicit function of that point and its prox image, cf.~\eqref{eq:right-bm-prox-identity}. Inspecting the proof of Lemma~\ref{lem:right-bm-locality}, the prox images coincide near~$x_s$, namely~$\prox_{\eta g_{T+1}}(x_s)=\prox_{\eta g_s}(x_s)$; therefore~$\nabla f(x_s)=\nabla f_s(x_s)$, consistently with the affine identity above. Since both the value and the gradient agree, the first-order oracle~\eqref{eq:oracle} returns identical data,
\begin{equation}
\label{eq:oracle-equal}
\Orc_{f_{T+1}}(x_s)=\Orc_{f_s}(x_s) \qquad \forall s \in [T].
\end{equation}
%
We argue by induction that, when run on~$f_{T+1}$,~$\Mtd$ queries exactly~$x_1,\dots,x_T$ and reports~$x_{T+1}$. For the base case, the first query~$x_1=\Phi_1$ is fixed by~$\Mtd$ independently of the objective, cf.~\eqref{eq:method-iterates}, so it coincides with the constructed one. 
Assume, for some~$2 \le t \le T$, that the first~$t-1$ queries~$x_1,\dots,x_{t-1}$ coincide with those in the transcript of~$\Mtd$ run on~$f_{T+1}$, and the corresponding constructed answers~$\Orc_{f_{1}}(x_1), \dots, \Orc_{f_{t-1}}(x_{t-1})$ coincide with those in the transcript,~$\Orc_{f_{T+1}}(x_1), \dots, \Orc_{f_{T+1}}(x_{t-1})$. 
By~\eqref{eq:method-iterates}, the next query of~$\Mtd$ given the constructed data is~$x_{t}=\Phi_{t}(\Orc_{f_{1}}(x_1),\dots,\Orc_{f_{t-1}}(x_{t-1}))$, 
and the next query in the transcript is~$\smash{\Phi_{t}(\Orc_{f_{T+1}}(x_1),\dots,\Orc_{f_{T+1}}(x_{t-1}))}$, identical by the induction premise. 
Now due to~\eqref{eq:oracle-equal}, the next constructed answer~$\smash{\Orc_{f_{t}}(x_{t})}$ is identical to the corresponding answer~$\smash{\Orc_{f_{T+1}}(x_{t})}$ in the transcript. 
This advances the induction. 
Finally, the same argument applied for~$t = T+1$, with the report map~$\smash{\bar\Phi_{T+1}}$ in place of~$\Phi_t$, gives consistency of the reported point. 
%and completes the proof.
%Applying the report map~$\bar\Phi_T$ to the matching answers at~$x_1,\dots,x_T$ yields the reported point~$x_{T+1}$, completing the proof.
%By~\eqref{eq:oracle-equal} the answer at~$x_t$ on~$f$ equals~$\Orc_{f_t}(x_t)$, i.e.~the constructed answer; the next query on~$f$ is then~$x_{t+1}=\Phi_{t+1}(\Orc_f(x_1),\dots,\Orc_f(x_t))$ by~\eqref{eq:method-iterates}, and since~$\Phi_{t+1}$ is a deterministic map applied to the same arguments, this equals the constructed~$x_{t+1}$. This advances the induction. 
\end{proof}

\subsection{Bounding the suboptimality gap}
\label{sec:gap-proof}

In this section, we complete the proof of Theorem~\ref{th:lower} by estimating the suboptimality gap of~$f_{T+1}$ at the reported point~$x_{T+1}$. 
To this end, we first bound~$f_{T+1}(x_{T+1})$ from below, by using Lemmas~\ref{lem:entry-growth}--\ref{lem:loss} and Proposition~\ref{prop:fresh-entries};
%We first control the smoothing loss through the largest active coordinate (Lemma~\ref{lem:loss}), 
then we exhibit a comparator whose support includes an unexplored direction.

%\begin{proof}[Proof of Theorem~\ref{th:lower}]
%By the homogeneity of the minimax risk recorded in Section~\ref{sec:formal}, it suffices to treat~$L=1$; let~$f\in\Fh(1)$ be the instance~\eqref{eq:final-instance}. Write~$x:=x_{T+1}$ and~$u:=\prox_{\eta g_{T+1}}(x)$.

\proofstep{1}:~{\em Lower bound at the reported point.} 
%Since~$g_{T+1}=\max_{s \in [T+1]} a_s \ge \ell_{T+1}$ and 
Recall that
$
\langle x_{T+1}, e_{j_{T+1}} \rangle \le \veps_T
$
by Proposition~\ref{prop:fresh-entries}, whence 
\begin{align}
\label{eq:final-atom}
	g_{T+1}(x_{T+1})
	\ge a_{T+1}(x_{T+1})
	&= 
	-\langle x_{T+1}, e_{j_{T+1}} \rangle - T\delta_T \notag\\
	&\ge
	-\veps_T - T\delta_T.
\end{align}
%Recall that~$a_s(\cdot)$ is \emph{active at}~$u$ when~$g_{T+1}(u)=a_s(u)$ (equivalently,~$s \in K_{g_{T+1}}(u)$ in the notation introduced after~\eqref{eq:backbone}): the corresponding~$j_s$ is then called an active coordinate. 
Let us bound each~$g_{T+1}$-{active coordinate} (of~$\hat u_{T+1}$) at the proximal point~$\hat u_{T+1}$ of~$x_{T+1}$, i.e.~$\langle \hat u_{T+1}, e_{j_t} \rangle$ for~$t \in [T+1]$ such that~$a_t(\hat u_{T+1}) = g_{T+1}(\hat u_{T+1})$. 
For such~$t$, the final atom is dominated at~$\hat u_{T+1}$: one has
$
a_t(\hat u_{T+1}) \ge a_{T+1}(\hat u_{T+1}),
$
that is,~$- \langle \hat u_{T+1}, e_{j_t} \rangle -(t-1)\delta_T \ge -\langle \hat u_{T+1}, e_{j_{T+1}} \rangle -T\delta_T$. 
Rearranging,
\begin{equation*}
        \langle \hat u_{T+1}, e_{j_t} \rangle 
        \le \langle \hat u_{T+1}, e_{j_{T+1}} \rangle + (T+1-t)\delta_T 
        \le \langle \hat u_{T+1}, e_{j_{T+1}} \rangle + T\delta_T.
\end{equation*}
%the last using~$t \ge 1$. 
Lemma~\ref{lem:entry-growth}, applied to~$g = g_{T+1}$ and~$x = x_{T+1}$, gives~$\langle \hat u_{T+1}, e_{j_{T+1}} \rangle \le 2 \langle x_{T+1}, e_{j_{T+1}} \rangle \le 2\veps_T$, whence
\begin{equation*}
        \langle \hat u_{T+1}, e_{j_t} \rangle  \le 2\veps_T + T\delta_T
\end{equation*}
for every~$g_{T+1}$-active coordinate~$j_t$ at~$\hat u_{T+1}$. 
As such, the premise of Lemma~\ref{lem:loss} holds for~$g = g_{T+1}$ and~$x = x_{T+1}$, with~$M = 2\veps_T + T\delta_T$.
Combining that lemma (cf.~\eqref{eq:loss-bound}) with~\eqref{eq:final-atom} results in
\begin{equation}
\label{eq:env-lower}
\begin{aligned}
%        f_{T+1}(x_{T+1})
%        \stackrel{\eqref{eq:final-instance}}{=}
%		=
\Env_{1/2}[g_{T+1}](x_{T+1})
        \stackrel{\eqref{eq:loss-bound}}{\ge}
g_{T+1}(x_{T+1}) - \frac{M}{2}
        \stackrel{\eqref{eq:final-atom}}{\ge}
         -2\veps_T - \frac{3T}{2} \delta_T.
%        &= -\veps_T - \tfrac{3}{2} T\delta_T.
\end{aligned}
\end{equation}

\proofstep{2}:~{\em Upper bound at the comparator}. 
%
Consider~$\bar x = \frac{1}{T+1} \sum_{t \in [T+1]} e_{j_{t}}$, i.e.~uniform on the selected coordinates; note that this includes the coordinate~$j_{T+1}$ revealed {\em after}~$\Mtd$'s reported point~$x_{T+1}$. 
Clearly,~$\bar x \in \Simp$. 
At~$\bar x$, each atom evaluates as~$a_t(\bar x) = -\langle \bar x, e_{j_t}\rangle-(t-1)\delta_T = -\tfrac1{T+1}-(t-1)\delta_T$. 
%Since~$i_s$ is itself a selected coordinate. 
Since~$(t-1)\delta_T \ge 0$,  we get
\begin{equation*}
        g_{T+1}(\bar x)
        =
        \max_{t \in [T+1]}  a_t(\bar x)
%        \max_{t \in [T+1]}  \left \{ -\frac1{T+1}-(s-1)\delta_T \right\}
        \le-\frac1{T+1}.
\end{equation*}
Since~$\bar x$ is feasible in~\eqref{eq:envelope}, this bound extends to the envelope:
$
        \Env_{1/2}[g_{T+1}](\bar x) \le g_{T+1}(\bar x) = -\frac1{T+1}.
$
%since~$\Div(q,q)=0$; hence~$\min_{\Simp}\Env_\eta[g_{T+1}]\le-1/(T+1)$. 
Combining this with~\eqref{eq:env-lower}, we get 
\begin{align*}
\Env_{1/2}[g_{T+1}](x_{T+1})-\min_{x \in \Simp}\Env_{1/2}[g_{T+1}](x)
&\ge
\frac1{T+1}-2\veps_T-\frac{3T}{2}\delta_T.
\end{align*}
Plugging in the parameter values~$ \veps_T :=\frac{1}{8(T+1)^2}$ and $ \delta_T = \frac{1}{3(T+1)^2} $ from~\eqref{eq:params}, and using~\eqref{eq:final-instance}, we get 
%%~$2\veps=\tfrac1{4(T+1)^2}$ and~$\tfrac32 T\delta=\tfrac{T}{2(T+1)^2}$, so that~$2\veps+\tfrac32 T\delta=\tfrac{2T+1}{4(T+1)^2}$, and writing~$\tfrac1{T+1}=\tfrac{4(T+1)}{4(T+1)^2}$ to reduce over the common denominator~$4(T+1)^2$,
%\begin{equation}\label{eq:gap}
%\Env_\eta[g_{T+1}](x)-\min_{\Simp}\Env_\eta[g_{T+1}]
%\ge
%\frac{4(T+1)-(2T+1)}{4(T+1)^2}
%=
%\frac{2T+3}{4(T+1)^2}.
%\end{equation}
%Finally,~$f=\eta\,\Env_\eta[g_{T+1}]$ by~\eqref{eq:final-instance}, so multiplying~\eqref{eq:gap} by~$\eta=\tfrac12$,
\begin{equation*}
%\label{eq:final-gap}
\begin{aligned}
        f(x_{T+1})-\min_{x \in \Simp}f(x)
        \ge
        \frac{1}{2}\left( \frac1{T+1}-2\veps_T-\frac{3T}{2}\delta_T \right)
        &= 
         \frac{1}{2(T+1)}\left(1-\frac{1}{4(T+1)}-\frac{T}{2(T+1)} \right) \\
        &=
         \frac{2T+3}{8(T+1)^2}
         > \frac{1}{4(T+1)}.
\end{aligned}
\end{equation*}
%the last step because~$2T+3\ge2(T+1)$. 
This completes the proof in the case of~$L = 1$. Recall that the general case follows by homogeneity.
%We complete the proof by reminding that the general case~$L > 0$ is recovered from 
%By Proposition~\ref{prop:transcript}, the left-hand side of~\eqref{eq:final-gap} is the genuine optimality gap of~$\Mtd$ on~$f\in\Fh(1)$; this proves~\eqref{eq:lower} for~$L=1$ and, by the homogeneity of the minimax risk recorded in Section~\ref{sec:formal}, for all~$L>0$. 
%
%We have thus established Theorem~\ref{th:lower}, and with it Theorem~\ref{th:intro-lower}.
\qed
%\end{proof}

\section{Quantum extension}
\label{sec:quantum}

%Let us now state this result rigorously and sketch its proof.

Let~$\Herm$ be the vector space of Hermitian~$d \times d$ matrices with inner product~$\inner{U}{V}:=\tr(UV)$, and let~$\Herm_+$ be the positive-semidefinite cone in~$\Herm$. 
Define the spectrahedron of density matrices
\begin{equation}
\label{eq:dens-def}
\Dens:=\{X\in\Herm_+:  \tr(X) = 1\}.
\end{equation}
The set~$\smash{\ri(\Dens)}$ consists of positive-definite unit-trace matrices in~$\smash{\Herm}$.
To formulate the result, it is convenient to use the notion of spectral functions.
Recall that any symmetric function~$\smash{\phi: \R^d \to \R}$ defines the spectral function~$\phi \circ \lambda$ on~$\Herm$,
where~$\lambda(X) = (\lambda_{1}(X),\dots,\lambda_d(X))$ is the ordered spectrum of~$X$.
It is well-known (e.g.~\cite[Thm.~1.1]{lewis1996derivatives}) that the gradient of~$\phi \circ \lambda$ at~$X = Q \Diag(\lambda(X)) Q^\dagger$ reads
\[
\nabla [\phi \circ \lambda](X) = Q \Diag(\nabla \phi(\lambda(X))) Q^\dagger
\] 
where~$A^\dagger$ is the conjugate transpose of a matrix~$A$ and~$\Diag(\cdot)$ maps~$x \in \R^d$ to the diagonal~$d \times d$ matrix with~$x$ on the diagonal; 
in other words,~$\smash{\nabla [\phi \circ \lambda](X)}$ is a Hermitian matrix with the same eigenbasis as~$\smash{X}$, whose spectrum is given by the gradient of~$\phi$ evaluated at the spectrum of~$X$.
In particular, negative entropy~\eqref{eq:intro-entropy} gives the negative von Neumann entropy~$h(\lambda(X))=\tr(X\log X)$. 
It is also well-known (e.g.~\cite{baes2007convexity}) that strict convexity of~$\phi \circ \lambda$ is equivalent to that of~$\phi$; this allows to extend Bregman divergences to~$\Herm$ via
$D_{\phi \hspace{1pt}\circ \lambda}(X,U) := \phi(\lambda(X)) - \phi(\lambda(U)) - \langle \nabla [\phi \circ \lambda](U), X-U \rangle$. 
In particular, {\em (Umegaki's) quantum relative entropy} between~$X \in \Dens$ and~$U \in \ri(\Dens)$, as given by
\begin{equation}
\label{eq:qre-def}
\qre{X}{U}:=\tr(X(\log X-\log U)),
\end{equation}
extends KL divergence in the above sense~\cite{hiai1991proper} and can be further extended to~$\Dens \times \Dens$ as in~\eqref{eq:intro-kl}.
%~$\mathbf{X} \times \ri(\mathbf{X})$ for some subset of~$$
As such, the appropriate generalization of~$\Fh(L)$ is the class~$\Qd(L)$ of convex functions~$F: \Dens \to \R$ continuously differentiable on~$\ri(\Dens)$ and satisfying the following inequalities:
\begin{equation}
\label{eq:quantum-sandwich}
0\le F(X)-F(U)-\inner{\bnabla F(U)}{X-U} 
\le L\qre{X}{U} \qquad \forall (X,U) \in \Dens \times \ri(\Dens).
\end{equation}
Here, as in~\eqref{eq:intro-sandwich}, we let~$\bnabla F$ be the tangent gradient of~$F$ on~$\ri(\Dens)$, so that~$\bnabla F(U) \in \spanof{I_d}^\perp$. 
%cf.~\eqref{eq:grad-tangent}.
%Since the orthogonal complement of the tangent space is~$\spanof{I}$, the gradient in~\eqref{eq:quantum-sandwich} is defined modulo~$\spanof{I}$. 

We can now state a noncommutative generalization of Theorem~\ref{th:intro-lower}.
%As in~\eqref{eq:oracle}, the corresponding first-order oracle is
%\begin{equation}
%\label{eq:quantum-oracle}
%\Orc_F^{\mathbf{H}}(X) 
%:= 
%(F(X),\bnabla F(X)).
%\end{equation}

%We are now ready to state the main result of this section.

\begin{theorem}
\label{th:quantum}
Let~$d,L,T$ be as in the premise of Theorem~\ref{th:lower}. 
For every deterministic method that makes~$T$ queries~$X_1, \dots, X_T \in \ri(\Dens)$ of the oracle~$X \mapsto (F(X),\bnabla F(X))$ and returns~$X_{T+1}\in\Dens$, there exists~$F(\cdot)\in\Qd(L)$ such that
\begin{equation}
\label{eq:quantum-bound}
F(X_{T+1})-\min_{X\in\Dens}F(X)>
\frac{L}{4(T+1)}.
\end{equation}
\end{theorem}

\begin{proof}[Proof sketch]
\proofstep{1}. First of all, we observe that the function~$F: \Dens \to \R$ given by~$F(X) = f(\diag(X))$, where~$f$ is the hard instance~$f \in \Ent(L)$ from Theorem~\ref{th:lower} and~$\diag(X) = (X_{11}, \dots, X_{dd})$, belongs to the class~$\Qd(L)$. 
Indeed, convexity of~$F$ follows from that of~$f$, since~$\diag(\cdot)$ is a linear mapping.
On the other hand, defining the diagonal truncation operator~$\Pin(H) := \Diag(\diag(H))$, we have
\begin{equation}
\label{eq:dpi}
\Div(\diag(X),\diag(U))
%\stackrel{\eqref{eq:qre-def}}{=}
= \qre{\Pin(X)}{\Pin(U)}
\le\qre{X}{U} \qquad \forall (X,U) \in \Dens \times \ri(\Dens).
\end{equation}
Here, the identity holds by~\eqref{eq:qre-def} and~\eqref{eq:intro-kl}; the final estimate is by Lindblad's data processing inequality~\cite{Lindblad1975,Uhlmann1977}.  
Meanwhile, by the chain rule~$\bnabla F(U)$ is a diagonal matrix with~$\bnabla f(\diag(U))$ on the main diagonal; as a result, we have the following in terms of~$x = \diag(X)$ and~$u = \diag(U)$,
\begin{align*}
F(X)-F(U)-\inner{\bnabla F(U)}{X-U} 
%= F(X)-F(U)-\inner{\bnabla F(U)}{\Pi(X)-\Pi(U)} \\
=
f(x)-f(u)-\inner{\bnabla f(u)}{x-u}  
%&\overset{\eqref{eq:intro-sandwich}}{\le}
\le L\Div(x,u)
%\overset{\eqref{eq:dpi}}{\le}
\le L\qre{X}{U},
\end{align*}
as required in~\eqref{eq:quantum-sandwich}. 
As such,~$F(\cdot) = f(\diag(\cdot))$ indeed belongs to~$\Qd(L)$ whenever~$f \in \Fh(L)$. 

\proofstep{2}. To complete the proof, we note that a method~$\tilde \Mtd$ with queries~$X_t \mapsto (F(X_t),\bnabla F(X_t))$, when run on~$F(\cdot) = f(\diag(\cdot))$, receives answers~$(f(x_t),\Diag(\bnabla f(x_t)))$ where~$x_t = \diag(X_t)$. 
Since the matrix~$\Diag(\bnabla f(x_t))$ contains the same information as~$\bnabla f(x_t)$, the sequence
%\[
%x_1,f(x_1),\bnabla f(x_1),\, \dots,\, x_T, f(x_T),\bnabla f(x_T), x_{T+1}
%\]
$(x_t,f(x_t),\bnabla f(x_t))_{t \in [T]}$ 
%of the transcript of~$\tilde\Mtd$ 
is the transcript of {\em some}~$\smash{\Mtd \in \FOM_d(T)}$ run on~$\smash{f \in \Fh(L)}$ and returning~$x_{T+1}$.
It remains to pick~$f$ as in Theorem~\ref{th:lower} and combine~\eqref{eq:intro-lower} with the fact that~$\smash{\min_{X \in \Dens} f(\diag(X)) = \min_{x \in \Simp} f(x)}$.
%we get~\eqref{eq:quantum-bound}.
\end{proof}
%
In Appendix~\ref{app:quantum-proof} we give an explicit minimax formulation of Theorem~\ref{th:quantum} analogous to Theorem~\ref{th:lower}, together with the full presentation of the emulation argument in step \proofstep{2} of the above proof sketch.
%with~$\omega_t = (f(x_t),\bnabla f(x_t))$
%To prove the theorem, it remains to make two observations:
%\begin{itemize}
%\item[(a)] any~$T$-step method~$\tilde \Mtd$ with access to the oracle~$X \mapsto (F(X),\bnabla F(X))$ can be converted into a method~$\Mtd \in \FOM_d(T)$ that queries the diagonals~$x_t = \diag(X_t)$ of the queries of~$\tilde \Mtd$, and only uses the diagonal part~$f(X),\bnabla f(X)$. 
%convert a 
%\item[(b)]
%\end{itemize}

%We next record a few useful identities that will be used in the reduction. Fix an orthonormal basis and define
%\begin{equation}\label{eq:diag-def}
%\diag(X):=(X_{11},\ldots,X_{dd}),
%\qquad
%\Diag(x):=\diag(x_1,\ldots,x_d),
%\qquad
%\Pin(X):=\Diag(\diag(X)).
%\end{equation}
%The map~$\diag$ sends~$\ri(\Dens)$ into~$\ri(\Simp)$ and maps~$\Dens$ onto~$\Simp$. Moreover,
%\begin{equation}\label{eq:diag-properties}
%\diag(\Diag(x))=x,
%\qquad
%\inner{\Diag(a)}{X}=\inner{a}{\diag(X)},
%\qquad x\in\Simp,\ a\in\R^d,\ X\in\Herm.
%\end{equation}
%The map~$\Pin$ is the dephasing channel. By~\eqref{eq:qre-def} and~\eqref{eq:intro-kl}, the quantum relative entropy of~$\Pin(X)$ and~$\Pin(U)$ equals the classical divergence of their diagonals. 
%
%Hence, Lindblad's data-processing inequality~\cite{Lindblad1975,Uhlmann1977} yields
%\begin{equation}
%\label{eq:dpi}
%\Div(\diag(X),\diag(U))
%=\qre{\Pin(X)}{\Pin(U)}
%\le\qre{X}{U}
%\qquad \forall (X,U)\in\Dens\times \ri(\Dens).
%\end{equation}
%\input{outro}
\section{Pointwise minimal interpolant}
\label{sec:outro}

%\paragraph{Improvement via }
%Theorem~\ref{th:intro-lower} leaves a logarithmic gap with the entropic mirror descent guarantee~\eqref{eq:intro-upper-KL}. We conjecture that this gap is due to the present resisting oracle construction. 
%Let us discuss a perspective approach for tightening the lower bound and potentially computing the exact value of~$\Risk_d(T,L)$. 
%Focusing on the case~$L=1$ (which is w.l.o.g. due to~\eqref{eq:risk-homogeneity}), 
A promising approach towards improving Theorem~\ref{th:intro-lower} would adapt the primal-dual estimate function (PDEF) framework of Florea and Nesterov~\cite{florea2025optimal}
%\footnote{In fact, this framework can be seen as an interpretation of the earlier result by Drori~\cite{Drori2017exact}, who established the minimax oracle complexity of the class of convex functions~$L$-smooth with respect to~$\|\cdot\|_2$.} 
to the class~$\Fh(L)$. 
%
This framework is built around the {\em optimal interpolating lower model}, defined as the pointwise-minimal function in the class of convex functions with~$L$-Lipschitz gradient (w.r.t.~$\|\cdot\|_2$), that interpolates the given data~$\smash{\{(x_s,f_s,\xi_s)\}_{s \in [t]}}$. The idea is that such a model can be updated incrementally: the next primal point~$x_{t+1}$ is selected by minimizing the current model~$\hat f_t$; a resisting oracle then selects an answer~$(f_{t+1}, \xi_{t+1})$ compatible with~$\hat f_t$, and the model is updated to incorporate the new point. 
Here, we make the first step towards implementing this program, by deriving an explicit form of the optimal interpolating lower model for the class~$\Fh(L)$, as a consequence of the results of Dragomir~\cite[Chap.~5]{DragomirThesis2021}. 
We note that the remaining step
%expected to be way more challenging, 
would be to construct a resisting oracle that gives the tightest final lower bound.

Given the data~$\cI=\{(x_i,f_i,\xi_i)\}_{i \in [N]}$ with~$(x_i,f_i,\xi_i) \in \ri(\Simp) \times \R \times \Tan$, we define the quantities
\begin{equation}
\label{eq:interpol-quantities}
r_i := f_i - \langle \xi_i, x_i \rangle,
\qquad
w_i := \exp({L^{-1}}{(r_i \ones + \xi_i)}), 
\qquad
W_{\cI} := \conv\{w_1, \dots, w_N\}.
\end{equation}
(In the sequel,~$\exp(\cdot)$,~$\log(\cdot)$ and division are entrywise.)
The next lemma, derived from~\cite[Thm.~5.11]{DragomirThesis2021} and proved in Appendix~\ref{app:interpol}, provides explicit interpolation conditions for the class~$\Fh(L)$.
%i.e.~inequalities over~$\cI$ that are equivalent to the existence of~$f \in \Fh(L)$ such that~$f_k = f(x_k)$ and~$\nabla f(x_k) = g_k$ for all~$k \in [K]$. 

\renewcommand{\ip}[2]{\langle #1, #2 \rangle}

\begin{lemma}
%[Interpolation conditions]
%[Simplex entropy interpolation]
\label{lem:entropy-interpolation}
The existence of~$f \in \Fh(L)$ that satisfies~$f(x_i) = f_i$ and~$\bnabla f(x_i) = \xi_i$ for all~$i \in [N]$ is equivalent to the following inequalities in terms of~\eqref{eq:interpol-quantities}:
\begin{equation}
\label{eq:entropy-interpolation}
% r_i-r_j\ge \log\ip{x_i}{\exp(g_j-g_i)},
% \qquad\text{equivalently}\qquad
 \ip{x_i}{w_j/w_i}\le 1 \qquad \forall i,j \in [N].
\end{equation}
\end{lemma}

%\begin{proof}
%\end{proof}

Using these interpolation conditions, we now derive the pointwise minimal interpolant in~$\Fh(L)$.

\begin{proposition}%[Optimal log--convex-hull interpolant]
\label{prop:log-hull}
Assume~$\cI = \{(x_i,f_i,\xi_i)\}_{i \in [N]}$ satisfies~\eqref{eq:entropy-interpolation}.
The function~$f_{\cI}: \Simp \to \R$ defined by
\[
 f_{\cI}(x):=L\max_{w\in W_{\cI}}\ip{x}{\log w}   %, \qquad x\in\Simp
\]
belongs to~$\Fh(L)$, interpolates~$\cI$, 
%satisfies~$f_{\cI}(x_i)=f_i$ and~$\bnabla f_{\cI}(x_i)=\xi_i$ for all~$i\in[N]$,
and is pointwise minimal among all functions with these properties.
\end{proposition}

\begin{proof}
\proofstep{1}. Let us verify that~$f_{\cI} \in \Fh(L)$. 
The function~$f_{\cI}$ is convex, and one has
\[
Lh(x)-f_{\cI}(x)=L\min_{w\in W_{\cI}} \sum_{k \in [d]} (x)_k \log\frac{(x)_k}{(w)_k}.
\]
The summand is jointly convex in~$((x)_k,(w)_k)$ as the perspective of a convex function~$-\log(\cdot)$, so~$Lh-f_{\cI}$ is convex. 
For~$x\in\ri(\Simp)$, strict concavity of~$w\mapsto\ip{x}{\log w}$ on~$W_{\cI}$ gives a unique maximizer~$w(x)$ that continuously depends on~$x$. Whence by  Danskin's theorem,~$f_{\cI} \in C^1(\ri(\Simp))$ with
$
\bnabla f_{\cI}(x)=L\ProjTan (\log w(x)),
$
and therefore~$f_{\cI}\in\Fh(L)$.

\proofstep{2}. 
We show that~$f_{\cI}$ interpolates~$\cI$. 
%By Lemma~\ref{lem:entropy-interpolation},~$\ip{x_i}{w_j/w_i}\le1$. 
For all~$w \in W_{\cI}$, we have~$\ip{x_i}{w/w_i} \le 1$  by~\eqref{eq:entropy-interpolation}, whence by concavity of~$\log(\cdot)$,
\[
 \ip{x_i}{\log(w/w_i)}\le\log\ip{x_i}{w/w_i} \le 0.
\]
Thus~$w(x_i) = w_i$, therefore
$f_{\cI}(x_i)=L\ip{x_i}{\log w_i}=f_i$
and
$
\bnabla f_{\cI}(x_i)=L\ProjTan(\log w_i)=\xi_i,
$
cf.~\eqref{eq:interpol-quantities}.

\proofstep{3}.
Finally, consider arbitrary~$F\in\Fh(L)$ that interpolates~$\cI$, i.e.~$F(x_i)=f_i$ and~$\bnabla F(x_i)=\xi_i$ for all~$i\in[N]$. 
Fix~$x\in\ri(\Simp)$, and set~$\xi:=\bnabla F(x)$ and~$r:=F(x)-\ip{\xi}{x}$. Applying Lemma~\ref{lem:entropy-interpolation} to the interpolation data~$\cI \cup \{(x,F(x),\xi)\}$ we get
\begin{equation}
\label{eq:interpol-condition-old-to-new}
 \exp({L^{-1}}{r_i})
 \ip{x}{\exp({L^{-1}}{(\xi_i-\xi)})}
 \le \exp({L^{-1}}{r})
 \qquad \forall i \in[N]
\end{equation}
where~$r = F-\langle\xi,x\rangle$.
As a result, for~$w=\sum_{i \in [N]} \lambda_i w_i\in W_{\cI}$ we get,  by using the concavity of~$\log(\cdot)$, 
\begin{align*}
L\ip{x}{\log w}
 &\le \ip{\xi}{x}
 +L\log\Bigg(\sum_{i \in [N]} \lambda_i \exp({L^{-1}}{r_i})
 \ip{x}{\exp({L^{-1}}{(\xi_i-\xi)})}\Bigg)
\stackrel{\eqref{eq:interpol-condition-old-to-new}}{\le} \ip{\xi}{x}+r 
 = F(x).
\end{align*}
Maximizing over~$w$ we get~$f_{\cI}(x)\le F(x)$ on~$\ri(\Simp)$. 
For~$\bar x$ on the relative boundary of~$\Simp$, consider the segment~$x(t) = (1-t)\bar x + t y$ with~$y \in \ri(\Simp)$. Then by continuity of~$f_{\cI}$ and convexity of~$F$ we get~$f_{\cI}(\bar x)=\lim_{t \downarrow 0}f_{\cI}(x(t))\le\limsup_{t\downarrow0}F(x(t))\le \limsup_{t \downarrow 0} (1-t) F(\bar x) + t F(y) = F(\bar x)$.
\end{proof}
\section*{Acknowledgments}

J.~M.~Aguirre is supported by the NSF Graduate Research Fellowship under Grant No.~DGE-2039655.
D.~M.~Ostrovskii thanks Radu-Alexandru Dragomir, Adrien Taylor, and Alexandre d'Aspremont for interesting discussions pertaining to this problem.
%\input{extra}
\appendix
%\section{Emulating the full gradient of~$f$ extended to~$\R^d_+$}
\section{Objective extension to~$\R^d_+$ with full-gradient oracle access}
\label{app:conic}

\paragraph{Tangent gradients and~$1$-homogeneous extension over~$\R^d_+$.}
%To avoid engaging with the formalism of affine differential geometry, 
Recall that any function~$f$ defined on~$\Simp$ and differentiable in~$\ri(\Simp)$---including the hard instances presented in Section~\ref{sec:hard-subclass}---admits the~1-homogeneous extension on~$\R^d_+$ (see, e.g.,~\cite{ovcharov2018proper}), defined as follows: setting~$s(x) := \ip{\ones}{x}$,
%concretely, the extension of~$f_0 \in \Simp \to \R$ is
\begin{equation}
\label{eq:conic-extension}
f^+(x) = s(x) f (s(x)^{-1}x).
%+ \langle \ones, x \rangle - 1.
\end{equation}
% the nonnegative orthant~$\R^d_{+}$ and restricted to~$\Simp$. 
This function is differentiable in~$\R^d_{++}$. 
An explicit calculation gives its full gradient for~$x \in \ri(\Simp)$:
\begin{equation}
\label{eq:grad-conic}
\nabla_\star f^+(x) = \bnabla f(x) + (f(x) - \langle \bnabla f(x), x \rangle) \ones.
\end{equation}
%where the full gradient~$\nabla_* f(x)$ is defined ``modulo~$\ones$:'' adding to it any vector in~$\textup{span}\{\ones\}$ does not change~$\nabla_* f^+$. 
\iffalse
Projecting~$\nabla_\star f^+$ onto the tangent space~$\Tan := \{v \in \R^d: \langle \ones, v \rangle = 0\}$ gives the {\em tangent gradient} of~$f$.
% the affine hull of the simplex:
%and note that, 
%We only evaluate~$\nabla f$ and~$\bnabla f$ on~$\ri(\Simp)$.
%We shall only take gradients~
%In contrast to~$\nabla f(x)$, 
Note that tangent gradient~$\bnabla f(x)$ and the gradient~$\nabla f^+(x)$ of the~$1$-homogeneous extension, for~$x \in \ri(\Simp)$, are both defined by the values of~$f$ in a {relative} neighborhood of~$x \in \ri(\Simp)$.
%The latter formula shows that the oracle which gives 
%As the result, the oracle~$x \mapsto (f(x),\bnabla f(x))$ is {more powerful} than the oracle~$x \mapsto \nabla f^+(x)$.
%Equivalently,~$\bnabla f$ is 
%Note that~$\langle \nabla f(x)$  w \rangle$ for any~$x \in \ri(\Simp)$ and~$w \in \Tan$
%
%If we identify~$\Tan$ with~$\R^{d-1}$ via the standard homeomorphism,~$\bnabla f(x)$ becomes the interior gradient. 
%of the restriction of the extension to~$\Tan$;
%
%We shall be referring to~$\bnabla f$ as the {\em lateral gradient} of~$f$.
%affine subspace~$\{x \in \R^d: \langle \ones, x \rangle = 1\}$, i.e.~the affine hull of~$\Simp$. 
%
% instances will  restrictions on the simplex~$\Delta_d$
For this reason, we interact with~$f$ via {tangent gradient}~$\bnabla f$, as formalized in Section~\ref{sec:formal}.
%As we shall explain in Section~\ref{sec:formal}, this is without loss of generality
%
Moreover,~$\Orc_f(x)$ allows to compute the full gradient~$\nabla f_+(x)$ of the 1-homogeneous extension~$f_+$ of~$f$, cf.~\eqref{eq:conic-extension}. 
Indeed, subtracting~\eqref{eq:grad-tangent} from~\eqref{eq:grad-conic} and using that~$x - d^{-1}\ones \in \Tan$ for~$x \in \Simp$, we get
\begin{equation}
\label{eq:grad-extended}
\nabla f_+(x) 
= \bnabla f(x) + \big( f(x) - \langle \bnabla f(x), x - d^{-1} \ones \rangle \big) \ones.
%= \bnabla f(x) + f(x) \ones - \langle \bnabla f(x), x - d^{-1} \ones \rangle \ones.
\end{equation}
\fi
Thus, the oracle~$\Orc_f(x) = (f(x),\bnabla f(x))$ is at least as strong as the full gradient oracle~$x \mapsto \nabla_\star f^+(x)$ in the extended class~$\{\smash{f^+: f \in \Fh(L)}\}$, in the sense that one can emulate the latter by using the former's answer at the same query point.
Intuitively, this should imply that access to the full gradient of~$\smash{f^+}$ {\em cannot} lead to a faster convergence rate than access to the oracle~$\smash{\Orc_f}$, as~$\smash{\Orc_f}$ gives at least as much information as~$\nabla_\star f^+$.
%\footnote{In fact, for {\em this extended class} of objectives the two oracles are equivalent, as one can also express~$\Orc_f(x)$ via~$\nabla f_+(x)$.}
%This intuition turns out to be correct; let us now formalize it.
The result we shall present next formalizes this comparison and leverages Theorem~\ref{th:lower} to obtain a lower bound for the class of functions on the positive orthant, convex and smooth relative to the unnormalized negative entropy
$
\smash{h_\e(x) := \sum_{k \in [d]} (x)_k \log(x)_k-(x)_k}
$
on~$\R^d_+$. 
%for this implication. 
%the formal argument is given in Appendix~\ref{app:conic}.
%In Remark~\ref{rem:emulation}, we showed that the oracle~$\Orc_f$ for~$f \in \Fh(L)$ is stronger than the full gradient oracle~$x \mapsto \nabla f_+(x)$ for the~$1$-homogeneous extension~$f_+$ of~$f$, cf.~\eqref{eq:conic-extension}, in the sense that~$\nabla f_+(x)$ is a function of~$\Orc_f(x)$; cf.~\eqref{eq:grad-extended}. 
%In this section, we derive a straightforward implication of this fact in the context of optimization under entropy-relative smoothness, with either of these two oracles.
%To make this rigorous, we consider the following construction. 

%Define the unnormalized negative entropy on~$\R^d_{+}$ by
%Its full gradient and convex conjugate are, respectively,~$\nabla_{\star} h_\e(x)=\log x$ and $h_\e^*(\xi)= \langle \ones, \exp(\xi) \rangle$.
Define~$\smash{\Ent_{d}^+(L)}$ as the class of functions~$\smash{\tilde f: \R^d_+ \to \R}$ that are convex,~$\smash{C^1(\R^d_{++})}$, and such that 
\begin{equation}
\label{eq:sandwich-conic}
\tilde f(x)- \tilde f(u)-\inner{\nabla_\star \tilde f(u)}{x-u} \le LD_{h_\e}(x,u) \quad \forall (x,u) \in \R^d_+ \times  \R^d_{++}.
\end{equation}
Since~$\smash{\inner{\nabla_\star \tilde f(x)}{v} = \inner{\bnabla \tilde f(x)}{v}}$ for all tangent directions~$v \in \Tan$, and~$h_e(x) = h(x)-1$ for all~$x \in \Simp$, the restriction of~$\tilde f \in \smash{\Fh^+(L)}$ to~$\Simp$ belongs to~$\Fh(L)$. The next lemma gives a partial converse.
%
\begin{lemma}
\label{lem:conic-extension}
The 1-homogeneous extension $f^+(x)$ of any~$f \in \Fh(L)$, cf.~\eqref{eq:conic-extension}, belongs to~$\Fh^+(L)$.
\end{lemma}
%
\begin{proof}
$f^+$ is convex as the composition of the perspective transformation of~$f$ and a linear mapping. 
Clearly,~$f^+ \in C^1(\R^d_{++})$. 
Finally, by writing 
\[
Lh_\e(x)-f^+(x)
 =s(x) (Lh(y)-f(y))\big|_{y = s(x)^{-1}x} +Ls(x)\log s(x)-L s(x)
\]
we see that~$Lh_\e - f^+$ is convex, which is equivalent to~\eqref{eq:sandwich-conic}. 
Indeed, the first term is convex as the composition of the perspective transformation of~$Lh-f$ (which is convex) and a linear mapping.
\end{proof}
%Indeed, the former implies the latter by restricting to the affine subspace~$s(x) = 1$; conversely, if~$Lh - f$ is convex on~$\Simp$, then~$Lh_\e(x)-f^+(x)$ is convex by the above identity, where 
%
%Conversely,~$\smash{\Fh^+(L)}$ contains 1-homogeneous extensions of functions from~$\Fh(L)$, as per~\eqref{eq:conic-extension}. 
%Indeed, the extension~$f^+$ of~$f \in \Fh(L)$ is convex (as a composition of the perspective function with a linear map), coincides with~$f$ on~$\Simp$, and has the same tangent gradient on~$\ri(\Simp)$, cf.~\eqref{eq:grad-conic}, so the right-hand inequality in~\eqref{eq:intro-sandwich} gives~\eqref{eq:sandwich-conic}.


Similarly, we can extend the class of first-order methods, by considering collections of mappings
%comprised of (i) a rule for selecting the initial query~$x_0 \in \ri(\Simp)$; (ii) a collection of mappings
\begin{equation}
\label{eq:method-maps-conic}
\begin{aligned}
\Psi_t: \R^{d(t-1)}  &\to \ri(\Simp), \qquad t \in [T];\\
\bar  \Psi_{T+1}: \R^{dT} &\to \Simp,
\end{aligned}
\end{equation}
where~$\R^d$ replaces the oracle answer space~$\Omega = \R \times \Tan$ in~\eqref{eq:method-maps}; 
this corresponds to minimizing~$\tilde f$ over~$\Simp$.
Let~$\FOM_d^+(T)$ be the class of all methods defined by~\eqref{eq:method-maps-conic} and define the minimax risk
%consider the correspondingly generalized notion of minimax risk: 
\begin{equation}
\label{eq:minimax-risk-conic}
\Risk_{d}^+(T,L) \;\; 
:=
\inf_{\tilde\Mtd^{\vphantom{{{2^2}^2}}} \,\in\, \FOM_{d}^+(T)} 
\quad
\sup_{\tilde f \,\in\, \Fh^+(L)} 
\quad
\left\{ \tilde f(x_{T+1}^{\tilde \Mtd}(\tilde f))- \min_{x \in \Simp} \tilde f(x) \right\},
\end{equation}
where the iterate sequence is generated by running~$\tilde \Mtd$ with the full gradient oracle~$x \mapsto \nabla \tilde f(x)$, i.e.
\begin{equation}
\label{eq:method-iterates-conic}
x_1^{\tilde\Mtd} = \Psi_1^{\vphantom{\tilde\Mtd}}, 
\;\;
x_2^{\tilde\Mtd}(\tilde f) = \Psi_2^{\vphantom{\tilde\Mtd}}(\nabla \tilde f(x_1^{\tilde\Mtd})), 
\;\; \ldots, \;\; 
x_{T+1}^{\tilde\Mtd}(\tilde f) = \bar\Psi_{T+1}^{\vphantom{\tilde\Mtd}}(\,\nabla \tilde f(x_1^{\tilde\Mtd}),\; \dots,\; 
\nabla \tilde f(x_{T}^{\tilde\Mtd}(\tilde f))).
\end{equation}
%
%In these terms, we can show that the~$\nabla_\star$-oracle complexity of~$\Ent_d^+(L)$ cannot be smaller than the~$\Orc_f$-oracle complexity of~$\Fh(L)$. 
%In other words, we have the following  comparison:
%
\begin{proposition}
\label{prop:emulation}
The minimax risks defined in~\eqref{eq:minimax-risk-conic} and~\eqref{eq:minimax-risk} satisfy
$
\Risk_{d}^+(T,L) \ge \Risk_{d}(T,L).
$
\end{proposition}

\begin{proof}
Replacing~$\Ent_d^+(L)$ with the smaller (due to Lemma~\ref{lem:conic-extension}) class~$\{f^+: f \in \Fh(L)\}$, we get
\[
\begin{aligned}
\Risk_{d}^+(T,L) 
&\ge 
\inf_{\tilde\Mtd^{\vphantom{{{2^2}^2}}} \,\in\, \FOM_{d}^+(T)} 
\quad
\sup_{f \,\in\, \Fh(L) \vphantom{\wt \Fh(L)}} 
\quad
\left\{  f^+(x_{T+1}^{\tilde \Mtd}(f^+))- \min_{x \in \Simp} f^+(x) \right\} \\ 
&= 
\inf_{\tilde\Mtd^{\vphantom{{{2^2}^2}}} \,\in\, \FOM_{d}^+(T)} 
\quad
\sup_{f \,\in\, \Fh(L) \vphantom{\wt \Fh(L)}} 
\quad
\left\{  f(x_{T+1}^{\tilde \Mtd}(f^+))- \min_{x \in \Simp} f(x) \right\}
\ge \Risk_{d}(T,L).
\end{aligned}
\]
%Here,~$f^+$ denotes the conic extension of~$f$, and 
Here the identity holds since the mappings in~\eqref{eq:method-maps-conic} output points in~$\Simp$ (where~$f^+ = f$). 
For the last inequality, we use~\eqref{eq:grad-conic} and observe
%\[
%\nabla f_+(x) 
%= \bnabla f(x) + (f(x) - \langle \bnabla f(x), x - \tfrac{1}{d} \ones \rangle) \ones.
%\]
%~\eqref{eq:grad-conic}--\eqref{eq:grad-tangent} to express~$\nabla f^+(x)$ in terms of~$\Orc_f(x) = (f(x),\bnabla f(x))$:
%\begin{equation*}
%\nabla f^+(x) = \bnabla f(x) + (f(x) - \bnabla f(x)^\top x) \ones.
%\end{equation*}
that the execution on~$f^+$ of~$\smash{\tilde\Mtd \in \FOM_d^+(T)}$, as per~\eqref{eq:method-maps-conic},  produces the same query sequence
\iffalse
\[
\begin{aligned}
\tilde x_1 &= \Psi_1, \\
\tilde x_2 &= 
\Psi_2 \left(\bnabla f(\tilde x_1) + \big(f(\tilde x_1) -\langle \bnabla f(\tilde x_1),  \tilde x_1 - d^{-1} \ones \rangle \big) \ones \right), \\
%\;=\; \Phi_2^{\vphantom{\Mtd}}(\Orc_f(x_1^\Mtd)), \\
\tilde x_{T+1} &= \bar\Psi_{T+1}
\Big( \bnabla f(\tilde x_1) + \big(f(\tilde x_1) \;- \langle \bnabla f(\tilde x_1),  \tilde x_1 - d^{-1} \ones \rangle \big) \ones, \dots, \\
&\qquad \qquad \, \bnabla f(\tilde x_T) + \big(f(\tilde x_T) -\langle \bnabla f(\tilde x_T),  \tilde x_T - d^{-1} \ones \rangle \big) \ones \Big),
\end{aligned}
\]
\fi
as the execution on~$f$ of the method
$
\Mtd \in \FOM_d(T),
$
defined by
%whose mappings~$\Phi_t$ are inductively defined  by 
\newcommand{\bg}{\xi}
\[
\begin{aligned}
\Phi_1 &:= \Psi_1, \\
\Phi_2(v_1, \bg_1) 
&:= \Psi_2 \left(\bg_1 + \big(v_1 -\langle \bg_1,  \Phi_1 \rangle \big) \ones \right), \\
\Phi_3(v_1,\bg_1, v_2, \bg_2) 
&:= \Psi_3 
		\left(
			\bg_1 + \big(v_1 -\langle \bg_1,  \Phi_1 \rangle \big) \ones,  \;\;
			\bg_2 + \big(v_2 -\langle \bg_2,  \Phi_2(v_1,\bg_1) \rangle \big) \ones
		\right), \\
&\;\;\vdots \vspace{-0.1cm} \\ \vspace{-0.1cm}
\bar\Phi_{T+1}(v_1,\bg_1, \dots, v_T, \bg_T) \\
:= \bar\Psi_{T+1} 
		\big(
		\bg_1 + \big(v_1 - &\langle \bg_1, \Phi_1 \rangle \big) \ones, \;\dots,\;
		\bg_T + \big(v_T -\langle \bg_T,  \Phi_T(v_1,\bg_1, \dots, v_{T-1},\bg_{T-1}) \rangle \big) \ones
		\big).
\end{aligned}
\]
%whose~$\Phi_t$ is defined as the composition of~$\Psi_t$ and 
%depends on~$\nabla f$ only through~$\bnabla f$.)
This can be verified by induction over~$t$, in the same way as in the proof of Proposition~\ref{prop:transcript}.
This fact implies the desired inequality, since in~$\Risk_{d}(T,L)$ the infimum is over {\em all} methods in~$\FOM_d(T)$. 
\end{proof}
\section{Explicit formulation and emulation argument for Theorem~\ref{th:quantum}}
\label{app:quantum-proof}

We first adapt the definition of first-order methods to make them in line with the class~$\Qd(L)$, cf.~\eqref{eq:quantum-sandwich}. 
Namely, we define the methods in~$\smash{\FOM_d^{\mathsf{H}}(T)}$  as collections of mappings (cf.~\eqref{eq:method-maps}):
\begin{equation}
\label{eq:method-maps-quantum}
\begin{aligned}
\boldsymbol{\Phi_t}: \boldsymbol{\Omega}^{t-1}  &\to \ri(\Dens), \qquad t \in [T]; \\
\boldsymbol{\bar \Phi_{T+1}}: \boldsymbol{\Omega}^T &\to \Dens,
\end{aligned}
\end{equation}
where~$\bOmega = \R \times \spanof{I_d}^\perp$ and~$\bOmega^0$ is a singleton.  
A given method in~$\FOM_d^{\mathsf{H}}(T)$ sequentially queries
\begin{equation}
\label{eq:oracle-quantum}
\Orc_F^{\mathsf{H}}(X) 
:= 
(F(X),\bnabla F(X))
\end{equation}
and uses the responses to form the sequence (cf.~\eqref{eq:method-iterates})
\begin{equation}
\label{eq:method-iterates-quantum}
X_1^{\vphantom\Mtd} = \bPhi_1^{\vphantom\Mtd}, 
\;\;
X_2^{\vphantom \Mtd}(F) = \bPhi_2^{\vphantom\Mtd}(\Orc_F^{\mathsf{H}}(X_1^{\vphantom\Mtd})), 
\;\; \ldots, \;\; 
X_{T+1}^{\vphantom\Mtd}(F) = \bar\bPhi_{T+1}^{\vphantom\Mtd}(\,\Orc_F^{\mathsf{H}}(X_1^{\vphantom\Mtd}),\; \dots,\; \Orc_F^{\mathsf{H}}(X_{T}^{\vphantom \Mtd}(F))).
\end{equation}
%along with the responses~$\Orc_f(x_1^\Mtd),\Orc_f(x_2^\Mtd(f)),\dots,\Orc_f(x_{T}^\Mtd(f))$, is called {\em the transcript} of~$\Mtd$ on~$f$.
In these terms, Theorem~\ref{th:quantum} claims that~$\Risk_d(T,L)$, cf.~\eqref{eq:minimax-risk}, is a lower bound for the quantum risk
\begin{equation}
\label{eq:minimax-risk-quantum}
\Risk_{d}^{\mathsf{H}}(T,L) \; 
:=
\inf_{{\bPhi_1, \dots, \bPhi_T, \bar\bPhi_{T+1}}^{\vphantom{\mathbf{H}}}} 
\quad
\sup_{F \,\in\, \Fh^{\mathsf{H}}(L)} 
\quad
\left\{ F(X_{T+1}(F))- \min_{X \in \Dens} F(X) \right\},
\end{equation}
where the infimum is over~$\FOM_d^{\mathsf{H}}(T)$, cf.~\eqref{eq:method-maps-quantum}--\eqref{eq:method-iterates-quantum}. 
%To verify this lower bound, we first observe that 
By the argument in step~(\proofstep{1}) of the proof sketch of Theorem~\ref{th:quantum}, we have~$(f \circ \diag) \in \Fh^{\mathsf{H}}(L)$, whence
\begin{align}
\Risk_{d}^{\mathsf{H}}(T,L) \; 
&\ge 
\inf_{{\bPhi_1, \dots, \bPhi_T, \bar\bPhi_{T+1}}^{\vphantom{\mathsf{H}}}} 
\quad
\sup_{f \,\in\, \Fh(L)} 
\quad
\left\{ f(\diag(X_{T+1}(f \circ \diag)))- \min_{X \in \Dens} f(\diag(X)) \right\} \notag\\
&=
\inf_{{\bPhi_1, \dots, \bPhi_T, \bar\bPhi_{T+1}}^{\vphantom{\mathsf{H}}}} 
\quad
\sup_{f \,\in\, \Fh(L)} 
\quad
\left\{ f(\diag(X_{T+1}(f \circ \diag)))- \min_{x \in \Simp} f(x) \right\}\,.
\label{eq:quantum-risk-chain}
\end{align}
Combining~$\bnabla [f \circ \diag](X) =  \Diag(\bnabla f(\diag(X)))$ with~\eqref{eq:oracle-quantum}--\eqref{eq:method-iterates-quantum}, for~$X_t = X_t(f \circ \diag)$  we get
\[
\begin{aligned}
X_{t} 
%&= \bPhi_t(\Orc_F^\mathbf{H}(X_1),\dots,\Orc_F^\mathbf{H}(X_{t-1})) \\
&= \bPhi_t(f(\diag(X_1)), \Diag(\bnabla f(\diag(X_1))), \dots, f(\diag(X_{t-1})), \Diag(\bnabla f(\diag(X_{t-1}))))
\end{aligned}
\]
when~$2 \le t \le T$, and
\[
\begin{aligned}
X_{T+1} &= \bar\bPhi_{T+1}(f(\diag(X_1)), \Diag(\bnabla f(\diag(X_1))), \dots, f(\diag(X_{T})), \Diag(\bnabla f(\diag(X_{T})))).
\end{aligned}
\]
Since in both cases~$X_t$ enters the right-hand side  only through the diagonal~$x_t := \diag(X_t)$, we get
\[
\begin{aligned}
x_{t}(f) 
%&= \bPhi_t(f(x_1), \Diag(\bnabla f(x_1)), \dots, f(x_{t-1}(f)), \Diag(\bnabla f(x_{t-1}(f)))) \\
&= \Phi_t(f(x_1), \bnabla f(x_1), \dots, f(x_{t-1}(f)), \bnabla f(x_{t-1}(f))), \quad 2 \le t \le T; \\
x_{T+1}(f) 
%&= \bPhi_{T+1}(f(x_1), \Diag(\bnabla f(x_1)), \dots, f(x_{T}(f)), \Diag(\bnabla f(x_{T}(f)))) \\
&= \bar\Phi_{T+1}(f(x_1), \bnabla f(x_1), \dots, f(x_{T}(f)), \bnabla f(x_{T}(f))),
\end{aligned}
\]
with~$\Phi_t: \Omega^{t-1} \to \ri(\Simp)$ defined by~$\Phi_t(v_1, g_1,\dots, v_{t-1}, g_{t-1}) = \bPhi_t(v_1, \Diag(g_1),\dots, v_{t-1}, \Diag(g_{t-1}))$; ditto for~$\bar\Phi_{T+1}$.  
These maps define a method in~$\FOM_d(T)$ with transcript~$(x_t,f(x_t),\bnabla f(x_t))_{t \in [T]}$ consistent with~$f \in \Fh(L)$;  plugging the reported point~$\diag(X_{T+1}(f \circ \diag)) = x_{T+1}(f)$ in~\eqref{eq:quantum-risk-chain},
%we get
\[
\Risk_{d}^{\mathsf{H}}(T,L) \; 
\ge 
\inf_{{\Phi_1, \dots, \Phi_T, \bar\Phi_{T+1}}^{\vphantom{\mathbf{H}}}} 
\quad
\sup_{f \,\in\, \Fh(L)} 
\quad
\left\{ f(x_{T+1}(f))- \min_{x \in \Simp} f(x) \right\} 
= \Risk_d(T,L). \qquad \qed
\]
%This completes the proof of Theorem~\ref{th:quantum}. 
%
%Some observations about \eqref{eq:diag-def}--\eqref{eq:dpi} are in order. 
%First, $\Diag$ is a right inverse of $\diag$, while $\Pin =\Diag \circ \diag$ is the associated dephasing retraction. Hence, for every $f:\Delta_{d}\to\mathbb{R}$ and $F:=f\circ\diag$, we have
%\[
%F\circ\Diag=f\,,\qquad F\circ\Pin=F.
%\]
%Moreover, $\diag$ sends $\ri(\Dens)$ into $\ri(\Delta_{d})$. 
%Thus, the diagonal pullback preserves feasible objective values and converts every admissible matrix-oracle query into an admissible simplex-oracle query. Second, the second identity in $\eqref{eq:diag-properties}$ shows that $\operatorname{Diag}$ is the adjoint of $\diag$. This relation is compatible with the gradient ambiguities on the two affine domains, since
%\[
%\Diag (a+\alpha\mathbf{1})=\operatorname{Diag}(a)+\alpha I.
%\]
%Hence a simplex gradient determines canonically a gradient class in $\mathbb{H}_{d}/\operatorname{span}{I}$. Finally, \eqref{eq:dpi} shows that the divergence controlling relative smoothness is nonincreasing under dephasing. The next lemma formalizes the resulting preservation of the objective class, the first-order oracle response, and the optimal value under the pullback $f\mapsto f\circ\diag$.
%
%\begin{lemma}\label{lem:lift}
%Let~$f\in\Fh(L)$ and set~$F:=f\circ\diag$. Then~$F\in\Qd(L)$ and
%\begin{equation}\label{eq:lift-identities}
%\nabla F(X)
%=\Diag(\bnabla f(\diag(X)))
%\quad \forall X \in \ri(\Dens),
%\qquad
%\min_{X \in \Dens} F(X)=\min_{x \in \Simp} f(x).
%\end{equation}
%\end{lemma}
%
%\begin{proof}
%Convexity and differentiability of~$F$ follow from the corresponding properties of~$f$ and the linearity of~$\diag$ in~\eqref{eq:diag-def}. The chain rule and the second identity in~\eqref{eq:diag-properties} give the gradient identity in~\eqref{eq:lift-identities}.
%
%Fix~$(X,U) \in \Dens \times \ri(\Dens)$ and set~$x:=\diag(X)$ and~$u:=\diag(U)$. In view of~\eqref{eq:intro-sandwich},~\eqref{eq:diag-properties}, and~\eqref{eq:dpi},
%\begin{align*}
%0
%&\le F(X)-F(U)-\inner{\bnabla F(U)}{X-U} \\
%&=f(x)-f(u)-\inner{\bnabla f(u)}{x-u} 
%\overset{\eqref{eq:intro-sandwich}}{\le}L\Div(x,u)
%\overset{\eqref{eq:dpi}}{\le}L\qre{X}{U}.
%\end{align*}
%This establishes~\eqref{eq:quantum-sandwich}. Finally, the first identity in~\eqref{eq:diag-properties} shows that~$\diag$ maps~$\Dens$ onto~$\Simp$, which proves the minimum identity in~\eqref{eq:lift-identities}.
%\end{proof}
%
%We are now ready to prove Theorem~\ref{th:quantum}.
%
%\begin{proof}
%Let~$\Mtd$ be a deterministic~$T$-query method for~\eqref{eq:quantum-oracle}. Construct a vector method~$\bar \Mtd$ that runs~$\Mtd$ internally. When~$\Mtd$ queries~$X_t \in \ri(\Dens)$, the method~$\bar \Mtd$ queries~\eqref{eq:oracle} at~$x_t:=\diag(X_t)\in\ri(\Simp)$ and supplies~$\Mtd$ with
%\[
%(f(x_t),\Diag(\bnabla f(x_t))).
%\]
%When~$\Mtd$ reports~$X_{T+1}$, the method~$\bar\Mtd$ reports~$x_{T+1}:=\diag(X_{T+1})$. Since~$\Mtd$ is deterministic and the maps in~\eqref{eq:diag-def} are fixed,~$\bar\Mtd$ is a deterministic~$T$-query method in~$\FOM_d(T)$. Thus, Theorem~\ref{th:lower} yields an~$f\in\Fh(L)$ such that
%\begin{equation}\label{eq:vector-gap}
%f(x_{T+1})-\min_{x\in\Simp}f(x)>
%\frac{L}{4(T+1)}.
%\end{equation}
%
%Set~$F:=f\circ\diag$. Lemma~\ref{lem:lift} gives~$F\in\Qd(L)$, and the first identity in~\eqref{eq:lift-identities} shows that every answer supplied to~$\Mtd$ is exactly~$\Orc_F(X_t)$. Hence, the determinism of~$\Mtd$ and the oracle recursion~\eqref{eq:method-iterates} imply that running~$\Mtd$ on~$F$ reproduces the simulated transcript and reports~$X_{T+1}$. Using the second identity in~\eqref{eq:lift-identities} and~\eqref{eq:vector-gap}, we obtain
%\begin{align*}
%F(X_{T+1})-\min_{X\in\Dens}F(X)
%=f(x_{T+1})-\min_{x\in\Simp}f(x) 
%\overset{\eqref{eq:vector-gap}}{>}
%\frac{L}{4(T+1)},
%\end{align*}
%which proves~\eqref{eq:quantum-bound}.
%\end{proof}


\section{Proof of Lemma~\ref{lem:entropy-interpolation}} 
%(interpolation conditions)
\label{app:interpol}

%Recall that by Lemma~\ref{lem:conic-extension},
The gradient and convex conjugate of~$h_\e$ are~$\nabla_{\star} h_\e(x)=\log x$ and $h_\e^*(\xi)= \langle \ones, \exp(\xi) \rangle$, respectively. 
%Set~$s(x) := \ip{\ones}{x}$. 
%The 1-homogeneous extension $f^+(x)=s(x) f(s(x)^{-1} x)$ is convex on~$\R^d_+$. 
%Besides,
%\[
%Lh_\e(x)-f^+(x)
% =s(x) (Lh(y)-f(y))\Big|_{y = s(x)^{-1}x} +Ls(x)\log s(x)-L s(x)
%\]
%shows that convexity of~$Lh_\e - f^+$ is equivalent to convexity of~$Lh-f$. 
%Indeed, the former implies the latter by restricting to the affine subspace~$s(x) = 1$; conversely, if~$Lh - f$ is convex on~$\Simp$, then~$Lh_\e(x)-f^+(x)$ is convex by the above identity, where the first term is convex as the composition of the perspective transformation of a convex function~$Lh-f$ and a linear mapping.
%which is convex because the first term is the perspective of the convex function~$Lh-f$, while the remaining terms are convex. Conversely, the restriction to~$\Simp$ of any convex function that is $L$-smooth relative to~$h_\e$ and interpolates the lifted data belongs to~$\Fh(L)$ and interpolates~$\cI$.

Lift~$\cI = \{(x_i,f_i,\xi_i)\}_{i \in [N]}$ to $\tilde\cI := \{(x_i,f_i, \tilde \xi_i)\}_{i \in [N]}$, where~$\tilde \xi_i:= \xi_i + r_i\ones$; cf.~\eqref{eq:interpol-quantities}. 
By~\eqref{eq:grad-conic}, we have~$(f(x_i),\nabla f(x_i))= (f_i,\xi_i)$ if and only if~$(f^+(x_i),\nabla_{\star} f^+(x_i))=(f_i,\tilde \xi_i)$. 
Meanwhile,~\cite[Thm.~5.11]{DragomirThesis2021} gives equivalent conditions of the existence of~$\tilde f \in \Fh^+(L)$ interpolating the data~$\tilde\cI$: for all~$i,j \in [N],$
\begin{equation}
\label{eq:interpol-conditions-unnormalized}
 A_{ij}:=f_i-f_j-\ip{\tilde \xi_j}{x_i-x_j}
 \ge L D_{h_\e^*}\big(\log x_i-{L^{-1}}{(\tilde \xi_i-\tilde \xi_j)},\log x_i\big).
\end{equation}
Since~$x_i,x_j\in\Simp$, we have~$A_{ij}=f_i-f_j-\ip{\xi_j}{x_i-x_j}$. Direct expansion of the right-hand side gives
\begin{align*}
D_{h_\e^*}(\log x_i+{L^{-1}}{(\tilde \xi_j-\tilde \xi_i)},\log x_i)
 &=\exp({L^{-1}}{(r_j-r_i)})
	\ip{x_i}{\exp({L^{-1}}{(\xi_j-\xi_i)})}
    -\ip{x_i}{L^{-1}(\tilde \xi_j-\tilde \xi_i)}-1\\
 &=L^{-1}A_{ij}
 	+\exp({L^{-1}}{(r_j-r_i)})
   	\ip{x_i}{\exp({L^{-1}}{(\xi_j-\xi_i})}-1.
\end{align*}
Thus~\eqref{eq:interpol-conditions-unnormalized} is equivalent to
$
\exp({L^{-1}}{(r_j-r_i)})
\ip{x_i}{\exp({L^{-1}}{(\xi_j-\xi_i})} \le 1,
$
that is to~\eqref{eq:entropy-interpolation}; cf.~\eqref{eq:interpol-quantities}.

Let us draw conclusions. 
If~$f \in \Fh(L)$ interpolates~$\cI$, then there exists a function in~$\Fh^+(L)$, namely the 1-homogeneous extension~$f^+$ of~$f$, that interpolates~$\tilde \cI$, so~\eqref{eq:entropy-interpolation} holds.
Conversely,~\eqref{eq:entropy-interpolation} implies the existence of~$\smash{\tilde f \in \Fh^+(L)}$ that interpolates~$\tilde \cI$, but then the restriction~$f$ of~$\smash{\tilde f}$ to~$\Simp$ interpolates~$\cI$: indeed,~$\tilde f(x_i) = f(x_i)$ since~$x_i \in \Simp$, and~$\bnabla f(x_i) = \ProjTan(\nabla_{\star} \tilde f(x_i)) = \ProjTan(\tilde \xi_i) = \xi_i$. 
\qed






\bibliographystyle{plain}
\bibliography{references}

\end{document}
